\providecommand{\zerJ}{\mathcal{J}}
\providecommand{\zerK}{\mathcal{K}}
\providecommand{\zersinc}{\operatorname{sinc}}
\providecommand{\zergE}{\gamma_{\mathrm{E}}}
\providecommand{\zerLL}{\mathcal{L}}
\providecommand{\zerconv}{\operatorname{conv}}
\providecommand{\zerdiag}{\operatorname{diag}}
\providecommand{\zerSym}{\mathfrak{D}}
\providecommand{\zerLog}{\operatorname{Log}}
\section{Hermitian couplings of completed \texorpdfstring{$L$}{L}-functions: real zeros, inverse problems and transport}\label{sec:cpl}

Completed Dirichlet $L$-functions coupled along a finite graph, or more generally by a Hermitian
matrix, have only real zeros as soon as the inputs have. This section determines the electrical and
edge laws of the displacement of zeros under such couplings and the couplings that minimize it for a
given conductance budget, solves the inverse problem (the zeros of one coupled function determine an
unknown list of characters, the labels and the coupling, by an exact decoder that is stable under
perturbation), and describes how infinitely many zeros are transported as the coupling varies, with
sharp-window estimates and local Cauchy laws for the logarithmic derivatives
(\cref{zer:sec:graph,zer:sec:transport}). The finite constructions are linear algebra and are marked
\stU; \cref{src:thm:dgrh} enters where the completed functions must be real canonical products, that is,
where their zeros must lie on the critical line. Inverse problems for single $L$-functions and their
products are in \cref{sec:chz}.

\subsection{Graph-coupled real-rooted functions}\label{zer:sec:graph}
\sources{graph_coupled_zeros_2026-11-26; grh_exact_extensions_2026-11-27; grh_path_rigidity_2026-11-27; grh_inverse_spectral_2026-11-27; coupled-dirichlet-inverse-spectrum; dirichlet-coupled-zeros-and-prime-variances}

This subsection is a statement about real-rooted polynomials and real canonical products; the
only arithmetic input is that \cref{src:thm:dgrh} makes the phase-normalized completed Dirichlet
$L$-functions such products. A weighted graph combines a family of real-rooted functions by a
rooted-forest sum of the functions and their first derivatives. For polynomials there is an
explicit symmetric matrix whose characteristic polynomial is this sum, the zeros lie to the left
of the critical points of the product in ordered matching, and the mean displacement is an
electrical invariant of the graph independent of the input zeros. The ingredients are classical:
the matrix-forest identity \cite{ChebotarevShamis2023Forest}, compressions representing critical
points \cite{Pereira2003Differentiators}, preservation of real zeros \cite{BorceaBranden2009LeeYang},
matrix-valued Herglotz theory \cite{Reiffenstein2022Interlacing} and rank-one inverse identities
\cite{SimonovWoracek2024Local}. Conversely, the zeros of the coupled function determine an unknown
list of characters of any parities and all principal minors of the coupling (grouped sums of them when
labels repeat): uniformly from a finite window for a known list, by an exact and Lipschitz-stable linear
decoder, and from the positive ordinates when the competitor misses them at only $o(T)$ distinct
ordinates below $T$ (\cref{zer:thm:window,zer:thm:decode,zer:thm:stabledecoder,zer:cor:blindall}).

\paragraph{The polynomial construction.}
Let $G$ be a connected undirected graph on $n\ge2$ labelled vertices with positive edge
conductances $w_e$ (parallel edges combined by adding conductances), Laplacian $L_G$, Moore--Penrose
inverse $L_G^+$, spanning-tree sum $\tau(G)=\sum_T\prod_{e\in T}w_e$ and effective resistances
$R_{ij}=(e_i-e_j)^{\mathsf T}L_G^+(e_i-e_j)$. At vertex $i$ place a monic real-rooted polynomial
$p_i$ of degree $d_i\ge1$, and put $N=d_1+\dots+d_n$; repeated roots and roots shared by different
vertices are allowed. A rooted spanning forest is a spanning forest with one chosen root in each
component; write $\mathcal R(F)$ for its roots. Put
\begin{equation}\label{zer:eq:forestpoly}
 \mathcal P_G(z)=\sum_{F}\Bigl(\prod_{e\in F}w_e\Bigr)\Bigl(\prod_{i\in\mathcal R(F)}p_i'(z)\Bigr)
 \Bigl(\prod_{i\notin\mathcal R(F)}p_i(z)\Bigr),
\end{equation}
the sum over rooted spanning forests, the empty edge set included. List the input roots with
multiplicity, in vertex blocks, as $r_1,\dots,r_N$, and put $\Lambda=\zerdiag(r_1,\dots,r_N)$,
$V_{ai}=1$ if the root $a$ belongs to $p_i$ and $0$ otherwise ($V$ is $N\times n$), and
$Q=I_N-\mathbf 1_N\mathbf 1_N^{\mathsf T}/N$. On $\mathcal V=\mathbf 1_N^\perp$ define
\[
 A_\infty=(Q\Lambda Q)|_{\mathcal V},\qquad B_G=(QVL_G^+V^{\mathsf T}Q)|_{\mathcal V},\qquad K_G=A_\infty-B_G .
\]

\begin{theorem}[Compression and electrical transport; \stU]\label{zer:thm:electrical}
For every complex $z$,
\begin{equation}\label{zer:eq:compression}
 \mathcal P_G(z)=N\tau(G)\det_{\mathcal V}(zI-K_G).
\end{equation}
In particular $\mathcal P_G$ has exactly $N-1$ zeros, all real. Let
$\lambda_1\le\dots\le\lambda_{N-1}$ be these zeros and $\beta_1\le\dots\le\beta_{N-1}$ the zeros of
$(\prod_ip_i)'$. Then $0\le\beta_k-\lambda_k\le\norm{B_G}$ for every $k$, and
\begin{equation}\label{zer:eq:resistance}
 \sum_{k=1}^{N-1}(\beta_k-\lambda_k)=\frac1N\sum_{i<j}d_id_jR_{ij}.
\end{equation}
Consequently the empirical measures $\mu_G$ of the $\lambda_k$ and $\nu$ of the $\beta_k$ satisfy
\[
 W_1(\mu_G,\nu)=\frac{\sum_{i<j}d_id_jR_{ij}}{N(N-1)},
\]
a value independent of every input root, and $W_\infty(\mu_G,\nu)\le\norm{B_G}$, which is sharp over
the permitted inputs even with the graph and the degrees fixed. The mean of the zeros is
$N^{-1}\sum_ar_a-\sum_{i<j}d_id_jR_{ij}/(N(N-1))$.
\end{theorem}
\begin{proof}
First,
\begin{equation}\label{zer:eq:forestdet}
 \det\bigl(L_G+\zerdiag(x_1,\dots,x_n)\bigr)=\sum_F\Bigl(\prod_{e\in F}w_e\Bigr)\prod_{i\in\mathcal R(F)}x_i :
\end{equation}
add a vertex $0$ joined to each vertex $i$ by an edge of weight $x_i$; the matrix-tree theorem
identifies the left side with the spanning-tree polynomial of the enlarged graph, and removing $0$
from a spanning tree leaves a spanning forest each of whose components has exactly one vertex
formerly adjacent to $0$, its root. This is a weight-preserving bijection and a polynomial identity.
Away from the input roots substitute $x_i=p_i'/p_i$ and multiply by $p=\prod_ip_i$:
$\mathcal P_G=p\det(L_G+\zerdiag(p_i'/p_i))$. Since $V^{\mathsf T}(zI-\Lambda)^{-1}V=\zerdiag(p_i'/p_i)$,
the Schur complement gives
\[
 \mathcal P_G(z)=\det\begin{pmatrix}zI_N-\Lambda&V\\-V^{\mathsf T}&L_G\end{pmatrix},
\]
first off the roots and then everywhere. Choose an orthonormal basis of $\R^n$ beginning with
$u=\mathbf 1_n/\sqrt n$ and otherwise consisting of eigenvectors of $L_G$, and eliminate the block of
positive eigenvalues. This gives $\det{}'(L_G)\det\begin{pmatrix}zI_N-\Lambda+VL_G^+V^{\mathsf T}&Vu\\-u^{\mathsf T}V^{\mathsf T}&0\end{pmatrix}$,
where $\det{}'(L_G)=n\tau(G)$ is the product of the positive eigenvalues. Now $Vu=\mathbf 1_N/\sqrt n$,
so $\norm{Vu}^2=N/n$ and $(Vu)^\perp=\mathcal V$, and for any $N\times N$ matrix $C$ and nonzero real
$b$ one has $\det\begin{pmatrix}C&b\\-b^{\mathsf T}&0\end{pmatrix}=\norm b^2\det_{b^\perp}(P_{b^\perp}C|_{b^\perp})$, by
expansion in a basis ending with $b/\norm b$. This proves \eqref{zer:eq:compression}.

The matrix $K_G$ is real symmetric and $B_G\succeq0$. The compression identity for derivatives,
$\det_{\mathcal V}(zI-A_\infty)=p'(z)/N$, follows from the cofactor formula
$p(z)e^{\mathsf T}(zI-\Lambda)^{-1}e$ with $e=\mathbf 1_N/\sqrt N$. Min--max applied to
$A_\infty-\norm{B_G}I\preceq K_G\preceq A_\infty$ gives the displacement bounds. With
$d=(d_i)$ and $D=\zerdiag(d_i)$, $V^{\mathsf T}V=D$ and $V^{\mathsf T}\mathbf 1_N=d$, so
\[
 \sum_k(\beta_k-\lambda_k)=\tr B_G=\tr\Bigl[L_G^+\Bigl(D-\frac{dd^{\mathsf T}}N\Bigr)\Bigr]
 =\sum_id_i(L_G^+)_{ii}-\frac1Nd^{\mathsf T}L_G^+d=\frac1N\sum_{i<j}d_id_jR_{ij},
\]
by $R_{ij}=(L_G^+)_{ii}+(L_G^+)_{jj}-2(L_G^+)_{ij}$. Since $\lambda_k\le\beta_k$, the ordered matching
attains the difference of the means, which bounds every coupling from below; this is the $W_1$
formula, and the ordered matching also gives the $W_\infty$ bound. For sharpness take
$p_i=(z-c)^{d_i}$, so that $A_\infty=cI$ and the largest displacement is $\norm{B_G}$. The mean
follows from $\tr A_\infty=(N-1)N^{-1}\sum_ar_a$.
\end{proof}

For the path on three vertices with $p_1=z+1$, $p_2=z$, $p_3=z-1$ and both conductances $t$,
$\mathcal P_t(z)=1+4tz+t^2(3z^2-1)$, with zeros $(-2\pm\sqrt{1+3t^2})/(3t)$ increasing to
$\pm1/\sqrt3$, and $W_1=2/(3t)$; at $t=1$ the zero $-4/3$ lies to the left of all three input
roots, so the theorem does not assert containment in the hull of the inputs. Exact rational checks
of this and the following examples are in
\artifact{graph_coupled_zeros_2026-11-26/code/verify_graph.py} (output
\artifact{graph_coupled_zeros_2026-11-26/checks/graph_checks.txt}).

\begin{theorem}[One edge; \stU]\label{zer:thm:edgelaw}
For distinct vertices $a,b$ put $b_e=e_a-e_b$, $R_e=b_e^{\mathsf T}L_G^+b_e$, $q_e=QVL_G^+b_e$, and
let $G_h$ have Laplacian $L_G+hb_eb_e^{\mathsf T}$, $h\ge0$. Then
\[
 K_{G_h}=K_G+\frac h{1+hR_e}q_eq_e^{\mathsf T},
\]
so $\lambda_k(G)\le\lambda_k(G_h)\le\lambda_{k+1}(G)$ for $k<N-1$ and
$\lambda_{N-1}(G)\le\lambda_{N-1}(G_h)$, and
\[
 \sum_{k=1}^{N-1}\bigl(\lambda_k(G_h)-\lambda_k(G)\bigr)=\frac h{1+hR_e}\norm{q_e}^2,\qquad
 \norm{q_e}^2=b_e^{\mathsf T}L_G^+\Bigl(D-\frac{dd^{\mathsf T}}N\Bigr)L_G^+b_e ,
\]
which is strictly positive for $h>0$, independent of the input roots, and increases to
$\norm{q_e}^2/R_e$ as $h\to\infty$. At a simple eigenvalue with unit eigenvector $u_k(h)$,
$\dd\lambda_k(G_h)/\dd h=\abs{u_k(h)^{\mathsf T}q_e}^2/(1+hR_e)^2$.
\end{theorem}
\begin{proof}
On $\mathbf 1_n^\perp$, where $L_G$ is positive definite, the rank-one inverse identity gives
$(L_G+hb_eb_e^{\mathsf T})^+=L_G^+-\frac h{1+hR_e}L_G^+b_eb_e^{\mathsf T}L_G^+$, and both sides vanish on
$\mathbf 1_n$. Substitution gives the matrix identity; monotonicity and rank-one interlacing give the
inequalities, and the trace gives the sum. If $q_e=0$, then $x=L_G^+b_e\in\mathbf 1_n^\perp$ has $Vx$
constant, so all coordinates of $x$ are equal because every vertex block is nonempty, whence $x=0$,
contradicting $L_Gx=b_e$. Differentiating at a simple eigenvalue gives the velocity.
\end{proof}

Scaling every conductance by $t$ gives $K_t=A_\infty-t^{-1}B_G$: all ordered zeros are nondecreasing in
$t$, $0\le\beta_k-\lambda_k(t)\le\norm{B_G}/t$, and $W_1(\mu_{tG},\nu)=\sum_{i<j}d_id_jR_{ij}/(tN(N-1))$. As
$t\downarrow0$, exactly $n-1$ zeros tend to $-\infty$ and the other $N-n$ tend, with multiplicity and
in order, to the zeros of $\prod_ip_i'$: on $\mathcal V$ the kernel of $B_G$ is $\ker V^{\mathsf T}$, of
codimension $n-1$, min--max sends the corresponding eigenvalues to $-\infty$, the complementary Schur
determinant tends to the characteristic polynomial of the compression of $\Lambda$ to
$\ker V^{\mathsf T}$, and that compression has characteristic polynomial $\prod_ip_i'/d_i$ by the
derivative compression applied blockwise. For $p_1=z^2-1$, $p_2=z^2-4$, $p_3=z^2-9$ on the path with
conductance $t$, $\mathcal P_t(z)=2z(3t^2z^4+8tz^3+(4-28t^2)z^2-36tz+49t^2)$ and the mean ordered
displacement from the critical points $0,\pm\sqrt{7/3},\pm\sqrt7$ is $8/(15t)$; adding an edge of
conductance $h$ between the ends moves the five zeros by $4h/(1+2h)$ in total.

\paragraph{Optimal couplings.}
Which graph moves the zeros least? Fix $n\ge2$, degrees $d_1,\dots,d_n\ge1$ and a budget $W>0$,
put $d=(d_i)$, $N=\sum_id_i$, $H=\mathbf 1_n^\perp\subset\R^n$ and
\begin{equation}\label{zer:eq:Cdeg}
 C=\zerdiag(d_1,\dots,d_n)-\frac{dd^{\mathsf T}}N .
\end{equation}
An admissible graph has the $n$ labelled vertices, conductances $w_{ij}=w_{ji}\ge0$ whose positive
support is connected, and $\sum_{i<j}w_{ij}=W$; the topology is free, its Laplacian $L$ is positive
definite on $H$, and $\tr L=2W$. For $1\le p<\infty$ let $\mathcal E_p(G)$ be the supremum, over all
monic real-rooted $p_i$ of degree $d_i$, of
\[
 W_p(\mu_G,\nu)=\Bigl(\frac1{N-1}\sum_{k=1}^{N-1}\abs{\beta_k-\lambda_k}^p\Bigr)^{1/p},
\]
and let $\mathcal E_\infty(G)$ be the supremum of $\max_k(\beta_k-\lambda_k)$; in dimension one the
ordered matching is optimal for every $W_p$, so these are Wasserstein distances between the
empirical measures of \cref{zer:thm:electrical}. They are guarantees depending only on the degrees.

\begin{lemma}[The reduced displacement operator; \stU]\label{zer:lem:reducedB}
The operator $B_G$ of \cref{zer:thm:electrical} is positive semidefinite of rank $n-1$, and its
nonzero eigenvalues are the eigenvalues on $H$ of
\[
 T_L=C^{1/2}(L|_H)^{-1}C^{1/2}.
\]
\end{lemma}
\begin{proof}
$V^{\mathsf T}QV=C$, so with $X=QV(L^+)^{1/2}$ the equal nonzero spectra of $XX^{\mathsf T}$ and
$X^{\mathsf T}X$ identify the positive spectrum of $B_G$ with that of $T_L$. The identity
$y^{\mathsf T}Cy=N^{-1}\sum_{i<j}d_id_j(y_i-y_j)^2$ shows that $C$ is positive definite on $H$, which
gives the rank.
\end{proof}

\begin{lemma}[Convex displacement bound; \stU]\label{zer:lem:convex}
Let $A$ be real symmetric and $B\succeq0$ of the same size $m$, order eigenvalues increasingly, and
put $\Delta_k=\lambda_k(A)-\lambda_k(A-B)\ge0$. For every continuous convex
$\phi\colon[0,\infty)\to\R$,
\[
 \sum_{k=1}^m\phi(\Delta_k)\le\tr\phi(B),
\]
with equality for $A=cI$, and for every $A$ when $\phi(t)=t$.
\end{lemma}
\begin{proof}
This is a form of Lidskii majorization \cite[Theorem~2.1]{Argentati2008Bounds}. Suppose first that $A$
has simple spectrum. The discriminant of the characteristic polynomial of $A-tB$ is a polynomial in
$t$, nonzero at $0$, so only finitely many $t\in[0,1]$ have repeated eigenvalues; elsewhere
$-\lambda_k'(t)=u_k(t)^{\mathsf T}Bu_k(t)=b_k(t)\ge0$ for an orthonormal eigenbasis $u_k(t)$. The
ordered eigenvalues are Lipschitz, so $\Delta_k=\int_0^1b_k(t)\,\dd t$ and Jensen's inequality gives
$\phi(\Delta_k)\le\int_0^1\phi(b_k(t))\,\dd t$. In an eigenbasis of $B$ each $b_k(t)$ is a convex
combination of the eigenvalues of $B$, with weights the squared entries of an orthogonal matrix, so a
second application of Jensen's inequality, summed over $k$, gives $\sum_k\phi(b_k(t))\le\tr\phi(B)$.
Integrate, and approximate a general $A$ by matrices with simple spectrum. The equality cases are
read off from the spectra and from the trace.
\end{proof}

Applied with $\phi(t)=t^p$ to $A_\infty$ and $B_G$, and with min--max for $p=\infty$, the lemmas give
the exact envelopes
\begin{equation}\label{zer:eq:envelope}
 \mathcal E_p(G)^p=\frac{\tr_H(T_L^p)}{N-1}\quad(1\le p<\infty),\qquad\mathcal E_\infty(G)=\norm{T_L},
\end{equation}
attained when every input root is one common real number, as in the sharpness part of
\cref{zer:thm:electrical}. Arbitrarily small distinct perturbations of a common root give the same
suprema, by continuity of eigenvalues.

\begin{theorem}[Exact minimax couplings; \stU]\label{zer:thm:minimax}
For $1\le p<\infty$ put $r=p/(p+1)$ and $S_r=\tr(C^r)$, the zero eigenvalue contributing zero. The
unique optimal graph has Laplacian
\[
 L_p=\frac{2W}{S_r}C^r,\qquad
 \min_G\mathcal E_p(G)=\frac{S_r^{1+1/p}}{2W(N-1)^{1/p}} .
\]
For the largest displacement the unique optimum is
\[
 L_\infty=\frac{2W}{\tr C}C,\qquad w^{(\infty)}_{ij}=\frac{Wd_id_j}{\sum_{a<b}d_ad_b},\qquad
 \min_G\mathcal E_\infty(G)=\frac{\tr C}{2W}=\frac{N^2-\sum_id_i^2}{2WN}.
\]
Every optimizer is a complete graph with strictly positive conductances. If the input roots are
required to be distinct, or the budget is only $\sum_{i<j}w_{ij}\le W$, the suprema and the optimal
graphs are unchanged.
\end{theorem}

At $p=1$ the displacement does not depend on the input roots, by \eqref{zer:eq:resistance}. For
$p>1$ the common root is an attaining configuration, not claimed to be the only one. The
normalization by $N-1$ matters: the perturbation $B_G$ has rank $n-1$, while there are $N-1$ output
zeros.

\begin{proof}
Work on $H$ in an orthonormal eigenbasis of $C$, with eigenvalues $c_1,\dots,c_{n-1}>0$, and put
$\ell_i=L_{ii}$ in this basis, so that $\sum_i\ell_i=\tr L=2W$. For finite $p\ge1$, the spectral
Jensen inequality and Cauchy--Schwarz give
\[
 \tr(T_L^p)\ge\sum_i\bigl((T_L)_{ii}\bigr)^p=\sum_i\bigl(c_i(L^{-1})_{ii}\bigr)^p
 \ge\sum_i\Bigl(\frac{c_i}{\ell_i}\Bigr)^p\ge\frac{\bigl(\sum_ic_i^{p/(p+1)}\bigr)^{p+1}}{(2W)^p}.
\]
Here $(L^{-1})_{ii}\ell_i\ge1$ is Cauchy--Schwarz for $L^{1/2}e_i$ and $L^{-1/2}e_i$, with equality
only if $e_i$ is an eigenvector of $L$; the last step is the minimum of the strictly convex function
$\sum_ic_i^p\ell_i^{-p}$ under $\sum_i\ell_i=2W$, attained only at $\ell_i=2Wc_i^r/S_r$. Equality
throughout forces $L=(2W/S_r)C^r$. At $p=1$ the first inequality is an equality, but the second still
forces diagonality and the scalar minimum is still unique. For $p=\infty$ let $t=\norm{T_L}$; then
$T_L\preceq tI$ is equivalent to $C\preceq tL$, so taking traces $t\ge\tr C/(2W)$, and at equality
$tL-C\succeq0$ has trace zero, hence vanishes.

These lower bounds hold for every positive definite $L|_H$ of trace $2W$, without the sign
conditions of a Laplacian; it remains to show that the candidates are graphs. For $0<r<1$,
\[
 C^r=\frac{\sin(\pi r)}\pi\int_0^\infty t^{r-1}C(C+tI)^{-1}\,\dd t ,
\]
which annihilates $\mathbf 1_n$, and for $i\ne j$ the rank-one inverse formula applied to
$C+tI=\zerdiag(d_i+t)-dd^{\mathsf T}/N$ gives
$((C+tI)^{-1})_{ij}=\frac{[d_i/(t+d_i)][d_j/(t+d_j)]}{t\sum_\ell d_\ell/(t+d_\ell)}>0$, so that
\begin{equation}\label{zer:eq:edgeintegral}
 -(C^r)_{ij}=\frac{\sin(\pi r)}\pi\int_0^\infty
 \frac{t^{r-1}d_id_j}{(t+d_i)(t+d_j)\sum_\ell d_\ell/(t+d_\ell)}\,\dd t>0 ,
\end{equation}
the integrand being $O(t^{r-1})$ at $0$ and $O(t^{r-2})$ at infinity. Thus $C^r$ is the Laplacian of
a complete graph with positive conductances; at $r=1$ this is $C_{ij}=-d_id_j/N$. Scaling gives
trace $2W$, and \eqref{zer:eq:envelope} gives the minima. Scaling all conductances by $c>0$ divides
every objective by $c$, so a graph using less than the budget is never optimal.
\end{proof}

At $p=\infty$ the optimum equalizes the displacements of the extremal input: for $p_i=(z-c)^{d_i}$ the
coupled zeros are $c-\tr C/(2W)$ with multiplicity $n-1$ and $c$ with multiplicity $N-n$. The
mean-optimal graph distributes them unequally and reduces their sum. Its edges keep the order of the
degree products but compress their contrasts.

Fractional powers of graph Laplacians are classical objects, with general preservation theorems
\cite[Theorem~3.1]{BenziSimunec2022Rational}; \eqref{zer:eq:edgeintegral} gives an explicit positive
integral for every optimal edge.

\begin{corollary}[Edges of the square-root network; \stU]\label{zer:cor:sqrtedges}
With $\eta(t)=\sum_kd_k/(t+d_k)$, the mean-optimal conductances are
\[
 w^{(1)}_{ij}=\frac{2Wd_id_j}{\pi\tr\sqrt C}\int_0^\infty\frac{\dd t}{\sqrt t\,(t+d_i)(t+d_j)\eta(t)},
\]
and for distinct $i,j,k$ with $d_j>d_k$,
\[
 1<\frac{w^{(1)}_{ij}}{w^{(1)}_{ik}}<\frac{d_j}{d_k}=\frac{w^{(\infty)}_{ij}}{w^{(\infty)}_{ik}};
\]
both ratios are one if $d_j=d_k$.
\end{corollary}
\begin{proof}
The formula is \eqref{zer:eq:edgeintegral} at $r=\tfrac12$. Cancelling common factors, the ratio
$w^{(1)}_{ij}/w^{(1)}_{ik}$ is the average of $(d_j/d_k)(t+d_k)/(t+d_j)$, which lies strictly between
$1$ and $d_j/d_k$ for every $t>0$, against the probability measure with density proportional to
$t^{-1/2}d_id_k/((t+d_i)(t+d_k)\eta(t))$.
\end{proof}

\begin{corollary}[Dependence on the displacement cost; \stU]\label{zer:cor:exponents}
If $n\ge3$, two distinct exponents $p\in[1,\infty]$ have the same optimal graph exactly when
$d_1=\dots=d_n$, and then every edge has conductance $2W/(n(n-1))$. For $n=2$ every exponent gives
the single edge of conductance $W$.
\end{corollary}
\begin{proof}
Equal normalized powers of $C$ with different exponents force all positive eigenvalues of $C$ to
coincide, so $C=c(I_n-\mathbf 1_n\mathbf 1_n^{\mathsf T}/n)$; its off-diagonal entries make all
products $d_id_j$, $i\ne j$, equal, which for $n\ge3$ means equal degrees. The converse and the case
$n=2$ are immediate.
\end{proof}

\begin{proposition}[Stability of the optimal couplings; \stU]\label{zer:prop:stability}
Let $L$ be an admissible Laplacian and $\norm\cdot_{S_1}$ the nuclear norm.
\begin{enumerate}[label=\textup{(\roman*)}]
\item If $\mathcal E_\infty(L)\le(1+\eps)\tr C/(2W)$, then
$\norm{L-L_\infty}_{S_1}\le4W\eps/(1+\eps)$ and $\norm{L-L_\infty}\le2W\eps/(1+\eps)$.
\item Let $L_*=L_1$, $\Delta=L-L_*$ and $\delta=\mathcal E_1(L)/\mathcal E_1(L_*)-1$. Then
\[
 \delta=\frac1{2W}\tr(\Delta L^+\Delta),\qquad\norm{L-L_*}_{S_1}\le2W\sqrt\delta ,
\]
and for $n\ge3$ the exponent $\tfrac12$ cannot be improved in a local bound uniform over admissible
Laplacians, in the nuclear or in the Frobenius norm.
\end{enumerate}
\end{proposition}
\begin{proof}
All inverses act on $H$; put $b=2W$. (i) Let $t=\mathcal E_\infty(L)$, $t_*=\tr C/b$ and
$\eta=t/t_*-1\le\eps$. By the proof of \cref{zer:thm:minimax}, $\Delta_\infty=tL-C\succeq0$, and
$L-L_\infty=\Delta_\infty/t-(t-t_*)C/(tt_*)$ is a difference of two positive semidefinite matrices of
trace $b\eta/(1+\eta)$; the triangle inequality gives the nuclear bound, and both terms lie below
$b\eta/(1+\eta)$ times the identity, which gives the operator bound.

(ii) With $\alpha=b/\tr\sqrt C$ one has $C=L_*^2/\alpha^2$, and the resolvent identity
$L^{-1}=L_*^{-1}-L_*^{-1}\Delta L_*^{-1}+L_*^{-1}\Delta L^{-1}\Delta L_*^{-1}$, multiplied by $C$, gives on
taking traces, since $\tr\Delta=0$,
$\tr(CL^{-1})-\tr(CL_*^{-1})=\alpha^{-2}\tr(\Delta L^{-1}\Delta)$ and $\tr(CL_*^{-1})=b/\alpha^2$, so
\eqref{zer:eq:envelope} gives the identity for $\delta$. Equivalently, with $\sigma=\tr\sqrt C$,
$\tr(CL^{-1})-\sigma^2/b=(\sigma/b)^2\norm{L^{-1/2}\Delta}_{\mathrm F}^2$, and
$\norm{XY}_{S_1}\le\norm X_{\mathrm F}\norm Y_{\mathrm F}$ with $X=L^{1/2}$, $\norm{L^{1/2}}_{\mathrm F}^2=b$,
gives $\norm\Delta_{S_1}\le(b\sqrt b/\sigma)(\delta\sigma^2/b)^{1/2}=b\sqrt\delta$. For sharpness exchange a
small amount of conductance between two edges of the strictly positive optimal complete graph:
$L(t)=L_*+tK$ with $K\ne0$ symmetric, $K\mathbf 1_n=0$ and $\tr K=0$. Then
$\delta(t)=t^2\tr(KL_*^+K)/b+O(t^3)$ with a positive coefficient, while both norms of $L(t)-L_*$ are
proportional to $\abs t$.
\end{proof}

For $(d_1,d_2,d_3)=(1,1,2)$ and $W=1$ the positive eigenvalues of $C$ are $1$ and $\tfrac32$, with
eigenvectors $(1,-1,0)$ and $(1,1,-2)$, and with $s=(3/2)^r$ the optimal conductances are
$w_{12}=(3-s)/(3(1+s))$ and $w_{13}=w_{23}=2s/(3(1+s))$. At $p=\infty$ they are
$(\tfrac15,\tfrac25,\tfrac25)$ with optimal largest displacement $\tfrac54$; at $p=1$,
$s=\sqrt6/2$, they are $(-3+4\sqrt6/3,\,2-2\sqrt6/3,\,2-2\sqrt6/3)$ with optimal mean displacement
$(5+2\sqrt6)/12$. The graph optimal for the maximum has mean $\tfrac56$, and the graph optimal for the
mean has maximum $(3+\sqrt6)/4>\tfrac54$. For the attaining inputs $(z,z,z^2)$ and conductances
$(x,y,y)$, $\mathcal P_G(z)=2z+(4x+6y)z^2+4(2xy+y^2)z^3$, which at the maximum-optimal graph is
$\tfrac{32}{25}z(z+\tfrac54)^2$, against $(p_1p_2p_3)'=4z^3$; at the mean-optimal graph the
displacements are $0$, $(1+s)/2$ and $s(1+s)/2$. These identities, a noncollapsed
forest-compression example and the stability identity are checked symbolically, together with
numerical comparisons of $75$ optimal configurations against $750$ competing graphs, in
\artifact{grh_exact_extensions_2026-11-27/code/verify_graph_minimax.py} (record
\artifact{grh_exact_extensions_2026-11-27/results/graph_verification.json}). For
$(d_1,d_2,d_3)=(1,1,4)$ and $W=1$ the positive eigenvalues of $C$ are $1$ and $2$; the minimax network has
$w_{12}=\tfrac19$, $w_{13}=w_{23}=\tfrac49$, largest displacement $\tfrac32$ and total displacement $3$,
and the mean network has $w_{12}=(3-\sqrt2)/(3(1+\sqrt2))$, $w_{13}=w_{23}=2\sqrt2/(3(1+\sqrt2))$, largest
displacement $(2+\sqrt2)/2$ and total displacement $(3+2\sqrt2)/2$, to be divided by $N-1=5$ for the
mean \artifact{grh_path_rigidity_2026-11-27/checks/network_checks.json}.

\paragraph{Entire functions.}
Let $F_1,\dots,F_n$ be nonzero entire functions of the form
\begin{equation}\label{zer:eq:realcanonical}
 F_i(z)=c_iz^{m_i}e^{b_iz}\prod_{\rho\in\mathcal Z_i}\Bigl(1-\frac z\rho\Bigr)e^{z/\rho},\qquad
 c_i\in\R\setminus\set0,\ b_i\in\R,\ m_i\ge0,
\end{equation}
with $\mathcal Z_i$ a multiset of nonzero reals, $\sum_{\rho\in\mathcal Z_i}\abs\rho^{-2}<\infty$, and
$m_i+\abs{\mathcal Z_i}>0$. This holds for real entire functions of order at most one with only real
zeros and at least one zero. Off the zeros
$g_i=F_i'/F_i=b_i+m_i/z+\sum_\rho(1/(z-\rho)+1/\rho)$, so that for $\Im z=y\ne0$
\begin{equation}\label{zer:eq:logderim}
 \Im g_i(z)=-y\Bigl(\frac{m_i}{\abs z^2}+\sum_{\rho\in\mathcal Z_i}\frac1{\abs{z-\rho}^2}\Bigr)
\end{equation}
has the sign of $-y$, and $-g_i'(x)=m_i/x^2+\sum_\rho(x-\rho)^{-2}>0$ at real regular points. The class
matters: $ze^{z^2}$ has only a real zero, while its derivative has nonreal zeros. For a constant
Hermitian $n\times n$ matrix $A$ put
\begin{equation}\label{zer:eq:PA}
 \mathcal P_A(z)=\det\bigl(\zerdiag(F_i(z))A+\zerdiag(F_i'(z))\bigr)
 =\sum_{S\subseteq[n]}\det A[S^c]\prod_{i\in S}F_i'(z)\prod_{j\notin S}F_j(z),
\end{equation}
a real entire function, and $M_A(z)=A+\zerdiag(g_i(z))$, so that
$\mathcal P_A=(\prod_iF_i)\det M_A$ off the input zeros.

\begin{theorem}[Hermitian coupling of real canonical products; \stU]\label{zer:thm:hermitian}
$\mathcal P_A$ is not identically zero and all its zeros are real; for $A=L_G$ it is the forest sum
\eqref{zer:eq:forestpoly} with $p_i$ replaced by $F_i$. For real $\lambda$ let
$I=\set{i:F_i(\lambda)=0}$, $J=[n]\setminus I$, $\mu_i=\ord_\lambda F_i$ for $i\in I$, and
$k_\lambda=\dim\ker\bigl(A[J]+\zerdiag(g_j(\lambda))_{j\in J}\bigr)$ (zero if $J$ is empty). Then
\begin{equation}\label{zer:eq:collision}
 \ord_\lambda\mathcal P_A=\sum_{i\in I}(\mu_i-1)+k_\lambda .
\end{equation}
If $a<b$ are not zeros of any $F_i$ or of $\mathcal P_A$, and $n_+$ counts positive eigenvalues,
\[
 N_{\mathcal P_A}(a,b)=\sum_iN_{F_i}(a,b)+n_+(M_A(a))-n_+(M_A(b)),
\]
where $N_F(a,b)$ counts zeros in $(a,b)$ with multiplicity; in particular the two counts differ by at
most $n$, and an interval without input zeros contains at most $n$ zeros of $\mathcal P_A$. Finally,
if $A(t)=A_0+tB$ with $B\succeq0$ and $x_0$ is a simple zero of $\mathcal P_{A(t_0)}$ with all
$F_i(x_0)\ne0$, the local real zero $x(t)$ satisfies
$x'(t_0)=v^*Bv/\sum_i\abs{v_i}^2(-g_i'(x_0))\ge0$ for any nonzero $v\in\ker M_{A(t_0)}(x_0)$.
\end{theorem}
\begin{proof}
For $\Im z>0$, if $M_A(z)v=0$ with $v\ne0$, then $0=\Im(v^*M_Av)=\sum_i\abs{v_i}^2\Im g_i(z)<0$. Hence
$\mathcal P_A\ne0$ in the upper half-plane, and in the lower one because it is real. The forest form
follows from \eqref{zer:eq:forestdet}.

For \eqref{zer:eq:collision} put $t=z-\lambda$ and $F_i(\lambda+t)=t^{\mu_i}h_i(t)$, $h_i(0)\ne0$, for
$i\in I$. The $I\times I$ block of $M_A(\lambda+t)$ is $t^{-1}D_I+E(t)$ with $D_I=\zerdiag(\mu_i)$ and $E$
analytic, so its inverse is $tD_I^{-1}+O(t^2)$ and its determinant $t^{-\abs I}\det D_I(1+O(t))$. The
Schur complement is $S(t)=C_\lambda+tD_\lambda+O(t^2)$ with
$C_\lambda=A[J]+\zerdiag(g_j(\lambda))$ and
$D_\lambda=\zerdiag(g_j'(\lambda))-A[J,I]D_I^{-1}A[I,J]\prec0$. On
$\ker C_\lambda\oplus(\ker C_\lambda)^\perp$ eliminate the invertible block; the remaining block is
$tD_\lambda|_{\ker C_\lambda}+O(t^2)$, negative definite to first order, so $\ord_0\det S=k_\lambda$. The
factors $\prod_{i\in I}F_i$ supply $\sum\mu_i$ and $\det$ of the $I$-block removes $\abs I$. The cases
$I=\varnothing$ and $J=\varnothing$ are the same argument with one block.

For the count, $M_A'(x)=\zerdiag(g_i'(x))\prec0$ off the input zeros, so at a zero of $\det M_A$ with
kernel of dimension $k$ exactly $k$ eigenvalues cross from positive to negative, and the positive
inertia drops by the multiplicity. Across an input zero, the $I$-block passes from $\abs I$ negative
to $\abs I$ positive eigenvalues, Hermitian elimination preserves inertia, and the $k_\lambda$ zero
eigenvalues of $S$ cross downwards, so $n_+$ jumps by $\abs I-k_\lambda=\sum_{i\in I}\mu_i-\ord_\lambda\mathcal P_A$.
Summing these jumps over the finitely many zeros in $[a,b]$ telescopes to the identity, and
$0\le n_+\le n$. For the velocity, the collision formula with $I=\varnothing$ shows the kernel is a
line; differentiating $M_{A(t)}(x(t))v(t)=0$ and multiplying on the left by $v^*$ gives
$v^*Bv+x'(t_0)\sum_i\abs{v_i}^2g_i'(x_0)=0$.
\end{proof}

The discrepancy bound is sharp for connected Laplacians: for $F_i(z)=z$ for all $i$,
$\mathcal P_{L_G}(z)=\prod_{j=2}^n(1+\ell_jz)$ with $\ell_j$ the nonzero eigenvalues of $L_G$, and an
interval close to $0$ contains $n$ input zeros and no output zero. If every input has a simple zero
at $\lambda$, then $\mathcal P_A(\lambda)\ne0$: the apparent common zero disappears.

\paragraph{Completed Dirichlet $L$-functions.}
Let $\chi$ be a primitive nonprincipal character of conductor $q$ and parity $a_\chi$, and
$\Lambda_\chi(s)=(q/\pi)^{(s+a_\chi)/2}\Gamma((s+a_\chi)/2)L(s,\chi)$, so that
$\Lambda_\chi(s)=\eps_\chi\Lambda_{\overline\chi}(1-s)$ with $\abs{\eps_\chi}=1$ and
$\overline{\Lambda_\chi(\overline s)}=\Lambda_{\overline\chi}(s)$ \cite{NISTDLMFDirichlet}. Choose
$\eta_\chi^2=\eps_\chi$ and put
\begin{equation}\label{zer:eq:phasereal}
 F_\chi(z)=\eta_\chi^{-1}\Lambda_\chi\bigl(\tfrac12+iz\bigr),\qquad
 F_\zeta(z)=\xi\bigl(\tfrac12+iz\bigr),\quad\xi(s)=\tfrac12s(s-1)\pi^{-s/2}\Gamma(s/2)\zeta(s).
\end{equation}
Then $\overline{F_\chi(\overline z)}=\eta_\chi\Lambda_{\overline\chi}(\tfrac12-iz)=\eta_\chi\eps_\chi^{-1}\Lambda_\chi(\tfrac12+iz)=F_\chi(z)$,
so $F_\chi$ is real entire even for complex $\chi$; the other square root changes only its sign. Under
\cref{src:thm:dgrh} all its zeros are real, and Hadamard factorization gives
\eqref{zer:eq:realcanonical} with real constants. The zero set is nonempty: a zero-free entire function
of order at most one is the exponential of a linear polynomial, whereas
$\log\abs{\Lambda_\chi(s)}=\tfrac s2\log s+O(s)$ on the positive real axis. With
$\eta_{\overline\chi}=\overline{\eta_\chi}$ one has $F_{\overline\chi}(z)=F_\chi(-z)$, so a conjugate
pair combines into one even input. Primitivity matters: imprimitive Euler factors have zeros on
$\Re s=0$, which are not real in the variable $z$.

\begin{corollary}[Graph-coupled completed $L$-functions; \stG]\label{zer:cor:graphdirichlet}
For any finite list of the functions \eqref{zer:eq:phasereal}, repetitions and $F_\zeta$ allowed, and any
Hermitian $A$, the entire function \eqref{zer:eq:PA} has only real zeros, its multiplicities are given
by \eqref{zer:eq:collision}, and on every interval with regular endpoints its zero count differs from
the sum of the input counts by at most $n$. For a graph Laplacian this is an explicit rooted-forest
combination of the completed functions and their first derivatives, all of whose zeros lie on
$\Re s=\tfrac12$ after $z=-i(s-\tfrac12)$. For every finite packet of zeros of the inputs, chosen with
their multiplicities, \cref{zer:thm:electrical,zer:thm:edgelaw} apply to the polynomials with these
roots, and the displacement budget depends only on the graph and the packet sizes.
\end{corollary}

In particular, for packets of $d_i$ zeros at vertex $i$ and a conductance budget $W$,
\cref{zer:thm:minimax} gives the best guarantee for every $W_p$ that depends only on the packet sizes
and the budget, whatever the multiplicities, coincidences and spacings of the selected zeros. A
particular arithmetic packet may have a smaller error and a different optimal graph, since the
common-root extremizer need not occur for Dirichlet $L$-functions; no optimum for an entire infinite
zero set is asserted.

No Euler product, automorphic interpretation or Riemann hypothesis for a further $L$-function is
asserted for these determinants, and central zeros need no exception, since their multiplicities
enter through $z^{m_i}$. For the odd characters $\chi_{-3},\chi_{-4}$ and the even character
$\chi_5$ the root numbers are $1$, so $F_3,F_4,F_5$ are even; on the path with conductances $a,b$,
\[
 \mathcal H_{a,b}=F_3'F_4'F_5'+aF_3F_4'F_5'+(a+b)F_3'F_4F_5'+bF_3'F_4'F_5+ab(F_3F_4F_5)'
\]
has only real zeros for all $a,b>0$. Grouped Dirichlet series show $F_q(0)>0$ for $q=3,4,5$ (for
$\chi_5$ each block $(5k+1)^{-s}-(5k+2)^{-s}-(5k+3)^{-s}+(5k+4)^{-s}$ is positive), so $F_q'(0)=0$, the
complementary matrix at $0$ is the path Laplacian of nullity one, and $0$ is a simple zero of
$\mathcal H_{a,b}$ by \eqref{zer:eq:collision}; directly,
$\mathcal H_{a,b}'(0)=abF_3(0)F_4(0)F_5(0)\sum_q(F_q'/F_q)'(0)<0$.

\paragraph{The inverse problem: characters and coupling from the zeros.}
\Cref{zer:cor:graphdirichlet} is a forward statement. Conversely, the zeros of $\mathcal P_A$ determine
the unknown list of characters and the principal minors of the coupling, grouped when labels repeat.
Parities may be mixed, the zeta input is allowed, a known list is read from a uniform finite window
after a bounded number of erasures, the coefficients are decoded by linear algebra even with false
observations, the decoder is Lipschitz stable under displacement of the observed ordinates, and an
unknown list is recovered from the positive ordinates even when the competitor misses the observed
zeros at $o(T)$ distinct ordinates below $T$.

Let $\mathcal X=(\chi_1,\dots,\chi_n)$ be a list of primitive characters, repetitions allowed, where
the character of conductor one stands for $F_\zeta$; let $F_i=F_{\chi_i}$ be as in
\eqref{zer:eq:phasereal} with arbitrary phases, $g_i=F_i'/F_i$, $Q_{\mathcal X}=\prod_iq_i$,
$S_{\mathcal X}(m)=\sum_i\chi_i(m)$, and for a Hermitian $A$ let $\mathcal P_A=\mathcal P_{\mathcal X,A}$ be
\eqref{zer:eq:PA} and $\nu_A(\gamma)=\ord_\gamma\mathcal P_A$. Changing a phase multiplies $\mathcal P_A$ by
$\pm1$. Let $\psi_1,\dots,\psi_h$ be the distinct labels, $m_j$ the number of vertices carrying $\psi_j$,
$m=(m_1,\dots,m_h)$, and $j(i)$ the label index of vertex $i$; inequalities between multi-indices are
componentwise. Put
\begin{equation}\label{zer:eq:groupD}
 D_A(X)=\det\bigl(A+\zerdiag(X_{j(1)},\dots,X_{j(n)})\bigr)=\sum_{0\le\alpha\le m}c_\alpha(A)X^\alpha,\qquad c_m=1,
\end{equation}
so that $c_\alpha(A)$ is the sum of the principal minors $\det A[J]$ over the $J$ whose complement has
$\alpha_\ell$ vertices of each label $\psi_\ell$, and
\[
 E_\alpha=\prod_{j=1}^h(F'_{\psi_j})^{\alpha_j}F_{\psi_j}^{m_j-\alpha_j},\quad
 \mathcal P_A=\sum_\alpha c_\alpha(A)E_\alpha,\quad
 V_\ell=\Bigl\{\sum_{\alpha\ne m}q_\alpha E_\alpha:\ q_\alpha\in\C[z],\ \deg q_\alpha\le\ell\Bigr\},
\]
with $\mathcal V=V_0$. For distinct labels $m=(1,\dots,1)$, the coefficient of $\prod_{i\notin J}X_i$ is
$\det A[J]$, and $\mathcal P_A=\sum_S\det A[S^c]E_S$ as in \eqref{zer:eq:PA}. If all labels coincide,
$D_A(X)=\det(A+XI)$, so even the exact zero set retains only the eigenvalues of $A$. For nonzero entire
$f,h$ and an interval $I$ let
\[
 \Delta_I(f,h)=\sum_{x\in I}\max\set{\ord_xf-\ord_xh,0},
\]
the one-sided count of multiplicities of real zeros of $f$ missing from $h$, possibly $+\infty$;
additional zeros of $h$ are not penalized.

\begin{lemma}[Independence on a vertical line; \stU]\label{zer:lem:algindep}
Let $\psi_1,\dots,\psi_h$ be distinct primitive characters, the conductor-one character and both parities
allowed, and $g_j=F_{\psi_j}'/F_{\psi_j}$. On the line $s=2+it$, $z=(s-\tfrac12)/i$, a polynomial
$P\in\C[X_1,\dots,X_h]$ with $P(g_1(z),\dots,g_h(z))\to0$ as $t\to+\infty$ is zero. In particular the
$g_j$ are algebraically independent over $\C(z)$, the functions $L'(s,\psi_j)/L(s,\psi_j)$ are
algebraically independent over $\C$, and the functions $E_\alpha$, $0\le\alpha\le m$, are linearly
independent over $\C(z)$.
\end{lemma}
\begin{proof}
Write $\Psi=\Gamma'/\Gamma$. The Euler product and the uniform digamma asymptotic give, for $t\ge1$,
\[
 g_j(z)=u(t)+c_j+D_j(t)+O((1+t)^{-1}),\qquad u(t)=\tfrac i2\log\frac{2+it}2,\quad c_j=\tfrac i2\log(q_j/\pi),
\]
with $D_j(t)=-i\sum_{k\ge2}\Lambda(k)\psi_j(k)k^{-2-it}$: the difference between $\Psi((s+a_j)/2)$ and
$\log(s/2)$ is $O(1/t)$ for either parity, and for $q_j=1$ the two extra logarithmic derivatives of
$s(s-1)$ also have that size. Each $D_j$ is a uniform limit of trigonometric polynomials by absolute
convergence.

For $d=\deg P$, replacing the $g_j$ by $u+c_j+D_j$ changes $P(g)$ by $o(1)$, since $\abs u=O(\log(2+t))$.
Expanding in $u$ gives $\sum_{k=0}^du^kf_k(t)=o(1)$, where the $f_k$ are bounded uniform limits of
trigonometric polynomials. Such a function with limit zero at $+\infty$ vanishes identically:
simultaneous recurrence of the finitely many phases of a uniform approximant gives arbitrarily large
translations arbitrarily close to periods. Since $\abs u\to\infty$, division by $u^d$ forces
$f_d(t)\to0$, hence $f_d=0$, and descending induction gives $f_0=P(c_1+D_1,\dots,c_h+D_h)=0$.

For $d>0$ let $\mathcal B$ polarize the top homogeneous term $P_d$ of $P$, with
$\mathcal B(x,\dots,x)=P_d(x)$, and extract the Fourier coefficient at $-\log k$ for $k=p_1\cdots p_d$ a
product of distinct primes; uniform absolute convergence permits termwise means. Each $D_j$ is
supported on prime powers, so only products from $P_d$ choosing one prime per factor contribute, and
the coefficient is
\[
 k^{-2}d!\,(-i)^d\prod_{r=1}^d\log p_r\;\mathcal B(v(p_1),\dots,v(p_d)),\qquad v(p)=(\psi_j(p))_{j=1}^h.
\]
The distinct primitive characters induce distinct characters on the units modulo
$Q=\lcm(q_1,\dots,q_h)$, so by orthogonality the vectors of character values span $\C^h$, and Dirichlet's
theorem supplies distinct primes realizing any prescribed tuple of unit classes, repeated classes
included. Multilinearity gives $\mathcal B=0$, a contradiction; a constant polynomial tending to zero is
zero. The same argument applied to a polynomial in the $c_j+D_j$ gives the independence of the
$L'/L(s,\psi_j)=iD_j$, after an affine change of variables.

Finally, clear denominators of a putative relation over $\C(z)$ and write it as
$\sum_{r=0}^bz^rP_r(g)=0$ with $P_b\ne0$. On the line $z=t-3i/2$, every lower term after division by $z^b$
is $O(t^{-1}(1+\log t)^{D})$ for some fixed $D$, so $P_b(g)\to0$, which is impossible. A linear relation
among the $E_\alpha$, divided by $\prod_iF_i$, is a relation among distinct monomials $\prod_jg_j^{\alpha_j}$.
\end{proof}

Ordinary linear independence of the $g_i$ would not suffice to recover determinants.

\begin{lemma}[Growth outside a strip; \stU]\label{zer:lem:outside}
Fix a Hermitian $A$, an integer $\ell\ge0$ and $H\in V_\ell$. Uniformly for $\abs{\Im z}\ge3/2$ and
$\abs z\to\infty$, $H(z)/\mathcal P_A(z)=O((1+\abs z)^\ell/\log\abs z)$. If $H\ne0$, the quotient
$H/\mathcal P_A$ is not a rational function.
\end{lemma}

\begin{proof}
For $\Im z\le-3/2$ one has $\Re s\ge2$; uniform Stirling asymptotics and the bounded logarithmic
derivative series there give $g_j(z)=\tfrac i2\log(s/2)+O(1)$ as $\abs s\to\infty$, for either parity.
All terms of $D_A$ below $X^m$ have total degree at most $n-1$, so its leading value is
$(\tfrac i2\log(s/2))^n(1+O(1/\log\abs s))$, which gives the bound; conjugation gives the same bound on
$\Im z\ge3/2$, conjugating the coefficients of $H$ when necessary. If $H/\mathcal P_A=r$ were rational,
dividing $H=r\mathcal P_A$ by $\prod_iF_i$ would give the relation
$\sum_{\alpha\ne m}q_\alpha g^\alpha=r\,D_A(g)$ over $\C(z)$, in which the coefficient of $g^m$ is $0$ on
the left and $r\ne0$ on the right, contrary to \cref{zer:lem:algindep}.
\end{proof}

\begin{lemma}[Finitely many poles force rationality; \stU]\label{zer:lem:invrational}
If $0\ne H\in V_\ell$ and $H/\mathcal P_A$ has only finitely many poles in $\C$, then $H/\mathcal P_A$ is
rational.
\end{lemma}

\begin{proof}
Let $W$ be the monic polynomial of degree $w$ cancelling the poles with their orders, and $J=WH/\mathcal P_A$.
The zero divisor of $\mathcal P_A$ is contained in that of $WH$, so by Hadamard factorization
\cite{Boas1954Entire} $J$ is a linear exponential times a canonical product of genus at most one over a
submultiset of the zeros of $WH$, whose counting function is $O_\eps(r^{1+\eps})$ by Jensen; thus $J$ is
entire of order at most one. By \cref{zer:lem:outside}, $J=O((1+\abs z)^{b})$ with $b=w+\ell$ outside the
strip $\abs{\Im z}<3/2$. Inside, for real $\beta>3/2$ the function $U=J/(z-i\beta)^{b}$ is bounded on the
boundary lines and of finite order on the strip; for $0<\kappa<\pi/3$ and $\delta>0$ the factor
$\exp(-\delta\cosh(\kappa z))$, whose exponent has real part at least $\cos(3\kappa/2)\cosh(\kappa\Re z)>0$
there, does not increase the boundary bound and makes the vertical sides of long rectangles tend to zero.
The maximum principle and $\delta\downarrow0$ give $J=O((1+\abs z)^{b})$ on $\C$, so $J$ is a polynomial by
Cauchy's estimates.
\end{proof}

\begin{theorem}[Tail divisor rigidity; \stG]\label{zer:thm:tailrigid}
For every Hermitian $A$, every $\ell\ge0$ and every nonzero $H\in V_\ell$, $\Delta_\R(\mathcal P_A,H)=+\infty$.
Hence, for Hermitian $A,B$ on the same list, $\Delta_\R(\mathcal P_A,\mathcal P_B)<\infty$ holds if and only
if $D_A=D_B$, that is $\mathcal P_A=\mathcal P_B$; for distinct labels this means
$\det A[J]=\det B[J]$ for every $J\subseteq[n]$.
\end{theorem}

\begin{proof}
By \cref{zer:cor:graphdirichlet} every zero of $\mathcal P_A$ is real. If $\Delta_\R(\mathcal P_A,H)<\infty$,
then $H/\mathcal P_A$ has finitely many poles and is rational by \cref{zer:lem:invrational}, contrary to
\cref{zer:lem:outside}. For $A,B$ put $H=\mathcal P_B-\mathcal P_A\in\mathcal V$; if $H\ne0$ then
$\ord_xH\ge\min(\ord_x\mathcal P_A,\ord_x\mathcal P_B)$, so $\Delta_\R(\mathcal P_A,H)\le\Delta_\R(\mathcal P_A,\mathcal P_B)<\infty$,
a contradiction. So $\mathcal P_A=\mathcal P_B$, and \cref{zer:lem:algindep} applied to the expansion in the
$E_\alpha$ gives $D_A=D_B$; the converse is trivial.
\end{proof}

\begin{theorem}[Uniform finite-window erasure correction; \stG]\label{zer:thm:window}
Fix the list $\mathcal X$. For every $M\ge0$ and integers $K,\ell\ge0$ there is a finite
$T_0=T_0(\mathcal X;M,K,\ell)$ such that $\Delta_{[-T_0,T_0]}(\mathcal P_A,H)>K$ for all Hermitian $A$ with
$\norm A\le M$ and all nonzero $H\in V_\ell$. In particular, if $B$ is any Hermitian matrix on the same
list, with no bound on its norm, and $\Delta_{[-T_0,T_0]}(\mathcal P_A,\mathcal P_B)\le K$ for the window with
$\ell=0$, then $D_A=D_B$.
\end{theorem}

\begin{proof}
Otherwise, for every $j$ there are a Hermitian $A_j$ with $\norm{A_j}\le M$ and $H_j\in V_\ell$ with unit
coefficient vector and $\Delta_{[-j,j]}(\mathcal P_{A_j},H_j)\le K$. By finite-dimensional compactness, along a
subsequence $A_j\to A_*$ and $H_j\to H_*$ locally uniformly; independence (\cref{zer:lem:algindep}) gives
$H_*\ne0$ and $\mathcal P_{A_*}\ne0$. By \cref{zer:thm:tailrigid} finitely many zeros of $\mathcal P_{A_*}$
carry excess over $H_*$ exceeding $K$. Take disjoint small discs about them, enclosing no other zero of
either limit, with boundaries free of zeros. Rouch\'e's theorem fixes the counts of $\mathcal P_{A_j}$ and
$H_j$ in these discs for all large $j$; every zero of $\mathcal P_{A_j}$ there is real and eventually in
$[-j,j]$, so the sum of the positive differences of the disc counts is a lower bound for
$\Delta_{[-j,j]}(\mathcal P_{A_j},H_j)$, a contradiction. Finally $\mathcal P_B-\mathcal P_A\in V_0$ vanishes to
at least the common multiplicity at every zero, so separation applied to this difference, if nonzero,
proves the last assertion.
\end{proof}

\begin{corollary}[Recovery of a labelled weighted graph; \stG]\label{zer:cor:graphrecover}
Let the labels be distinct. For real symmetric $A,B$ the conclusions of
\cref{zer:thm:tailrigid,zer:thm:window} hold if and only if $B=DAD$ for a diagonal $D$ with entries
$\pm1$. If $A,B$ are Laplacians of labelled graphs with nonnegative weights, then $A=B$: the one-vertex
minors give the weighted degrees and, for $i\ne j$, the weight is
$w_{ij}=\sqrt{L_{ii}L_{jj}-\det L[\set{i,j}]}$. Thus the weighted graph is determined by any cofinite
submultiset of the zeros of $\mathcal P_{L_G}$, or by a sufficiently long window after at most $K$ missing
multiplicities, with no bound on the competitor's weights.
\end{corollary}

\begin{proof}
Equal principal minors of order one and two give $a_{ii}=b_{ii}$ and $a_{ij}^2=b_{ij}^2$, so both
matrices have the same support graph, and $r_{ij}=b_{ij}/a_{ij}\in\set{\pm1}$ on its edges. On a chordless
cycle of length $\ell\ge3$ the principal minor on its vertex set is a sum over fixed points, disjoint edge
transpositions and the two orientations of the cycle; the first two kinds agree, and the cycle
contributes $2(-1)^{\ell-1}$ times the edge product, so $\prod_Cr_{ij}=1$, and by splitting at chords the
same holds on every simple cycle. Propagating signs $d_j=d_ir_{ij}$ along a spanning tree of each component
gives $B=DAD$; conversely a diagonal sign switch preserves all minors. Off-diagonal Laplacian entries are
nonpositive, so their magnitudes determine them. This is a case of the theorem of Engel and Schneider
\cite{EngelSchneider1980Matrices}.
\end{proof}

For signed symmetric matrices the switching ambiguity is unavoidable, since $D$ commutes with both
diagonal matrices in \eqref{zer:eq:PA}; for a general Hermitian matrix, diagonal unitary conjugation and
entrywise conjugation already preserve all principal minors, and the invariant recovered is their full
vector. A signed triangle with zero diagonal and unit edge magnitudes has
$\det(A+\zerdiag(X_i))=X_1X_2X_3-X_1-X_2-X_3+2\sigma$ with $\sigma$ the cycle sign, so the zeros detect cycle
signs but not vertex switches. For the even quadratic characters $\chi_5,\chi_8,\chi_{12}$ and the path
Laplacian $L(a,b)$ with weights $a,b$ on the edges $12,23$,
$\det(L(a,b)+\zerdiag(X_i))=X_1X_2X_3+aX_2X_3+(a+b)X_1X_3+bX_1X_2+ab(X_1+X_2+X_3)$; the Laplacians
$L(2,3)$ and $L(3,2)$ have the same ordinary spectrum, but their character-labelled coupled zero sets
differ, in each direction, at $\gg T$ distinct positive ordinates below $T$, by
\cref{zer:thm:localcoupled} below. The same holds for the mixed-parity labels
$\chi_{-3},\chi_{-4},\chi_5$ of $\mathcal H_{a,b}$, and by \cref{zer:cor:blindall} the zeros of
$\mathcal H_{a,b}$ identify all three characters, the middle one and both weights, against competing
triples of distinct characters of arbitrary conductors.

\paragraph{The prime frequencies of the coupled divisor.}
On the line $\Re s=\sigma>1$ put $b=\sigma-\tfrac12$, $P(s)=\mathcal P_A(-i(s-\tfrac12))$ and
\begin{equation}\label{zer:eq:poissonH}
 H_A(t)=\Re\frac{P'(\sigma+it)}{P(\sigma+it)}=\sum_\gamma\nu_A(\gamma)\frac{b}{b^2+(t-\gamma)^2},
\end{equation}
the sum over the distinct real zeros of $\mathcal P_A$: the real Hadamard product of $\mathcal P_A$ in $z$,
differentiated in $s$, gives \eqref{zer:eq:poissonH}, since the real exponential slope and each genus-one
correction become purely imaginary, and a zero at $z=0$ contributes the same Poisson term per copy.

\begin{theorem}[Landau formula for the coupled divisor; \stG]\label{zer:thm:coupledlandau}
For fixed real $x>1$ and $T\to\infty$,
\begin{equation}\label{zer:eq:coupledlandau}
 \frac1T\sum_{0<\gamma\le T}\nu_A(\gamma)x^{i\gamma}=-\frac{\Lambda(x)}{2\pi\sqrt x}S_{\mathcal X}(x)
 +O_{\mathcal X,A,x}\bigl((\log T)^{-1}\bigr),
\end{equation}
where the main term vanishes unless $x=p^k$ with $p$ prime and $k\ge1$, and then $\Lambda(x)=\log p$. With
$N_A(T)=\sum_{0<\gamma\le T}\nu_A(\gamma)$,
\begin{equation}\label{zer:eq:coupledcount}
 N_A(T)=\frac{nT}{2\pi}\Bigl(\log\frac T{2\pi}-1\Bigr)+\frac T{2\pi}\log Q_{\mathcal X}+O_{\mathcal X,A}(T/\log T).
\end{equation}
\end{theorem}
\begin{proof}
Take $\sigma=2$, $b=\tfrac32$. Stirling's formula for the digamma function gives, with
$U=\tfrac i2\log(s/2)$ and the principal logarithm, uniformly as $\abs t\to\infty$ on $s=2+it$,
\[
 g_j(-i(s-\tfrac12))=U+\tfrac i2\log(q_j/\pi)+iL_j'(s)/L_j(s)+O(1/\abs t),
\]
with $L_j=L(s,\chi_j)$; for the conductor-one input the extra term $i(1/s+1/(s-1))$ belongs to the error
and its derivative is $O(\abs t^{-2})$. These logarithmic derivatives and their first derivatives are
bounded on $\Re s=2$ by absolute convergence of their Euler series, and the differentiated gamma term is
$O(1/\abs t)$. The matrix $A+\zerdiag(g_j)$ is $UI+O(1)$, so its inverse has norm $O(1/\log\abs t)$, and
Jacobi's determinant identity bounds the $s$-logarithmic derivative of the determinant factor by
$O(1/\log\abs t)$. Thus, with $D_{\mathcal X}=\sum_jL_j'/L_j$,
\begin{equation}\label{zer:eq:cline}
 H_A(t)=\frac n2\log\frac{\abs t}{2\pi}+\frac12\log Q_{\mathcal X}+\Re D_{\mathcal X}(2+it)+O(1/\log\abs t),
\end{equation}
and \eqref{zer:eq:poissonH} with positivity gives $\sum_{u\le\gamma\le u+1}\nu_A(\gamma)\ll\log(2+\abs u)$
uniformly on $\R$. For fixed real $\omega$, using the Fourier transform of $b/(b^2+t^2)$,
\begin{equation}\label{zer:eq:unsmooth}
 \int_0^TH_A(t)e^{i\omega t}\,\dd t=\pi e^{-b\abs\omega}\sum_{0<\gamma\le T}\nu_A(\gamma)e^{i\omega\gamma}+O\bigl(\log^2(2+T)\bigr):
\end{equation}
for $-T\le\gamma\le2T$, replacing the integral for each zero by its full-line value times
$\mathbf 1_{(0,T]}(\gamma)$ costs $O((1+\abs\gamma)^{-1}+(1+\abs{T-\gamma})^{-1})$, which summed over unit
intervals costs $O(\log^2(2+T))$; outside this range the cost is $O(T/\gamma^2)$ and the local count makes
the tail $O(\log(2+T))$. This also justifies the exchanges.

At $\omega=\log x>0$ the logarithmic and constant terms of \eqref{zer:eq:cline} integrate to
$O_x(\log(2+T))$ by integration by parts. In the absolutely convergent series
$D_{\mathcal X}(s)=-\sum_{k\ge2}\Lambda(k)S_{\mathcal X}(k)k^{-s}$ the frequency $-\log x$ has coefficient
$-\Lambda(x)S_{\mathcal X}(x)x^{-2}$ if $x$ is a prime power, halved by taking the real part; every other
frequency contributes $O_x(1)$ in total after integration, since the nonzero numbers $\abs{\log(x/k)}$
for integers $k\ne x$ are bounded away from zero at fixed $x$, and the conjugate frequencies have the same
property. The integrated remainder is $O(T/\log T)$. Dividing \eqref{zer:eq:unsmooth} by $\pi Tx^{-b}$
proves \eqref{zer:eq:coupledlandau}. At $\omega=0$ the series has zero mean with integrated error $O(1)$,
and integration of the logarithm gives \eqref{zer:eq:coupledcount}.
\end{proof}

The leading term does not depend on the coupling: every prime frequency of the coupled zeros retains
the sum of the character values. This transfers the classical formula of Landau
\cite{Landau1911Nullstellen,KaczorowskiLanguascoPerelli2000Note} to the determinant, which generally has no
known Euler product. For two lists $\mathcal X,\mathcal Y$ and couplings $A,B$ write $X=Z_{\mathcal X,A}$ and
$Y=Z_{\mathcal Y,B}$ for the zero multisets and
\[
 M_{X\to Y}(T)=\sum_{0<\gamma\le T}\bigl(\operatorname{mult}_X(\gamma)-\operatorname{mult}_Y(\gamma)\bigr)_+,
\]
the sum over distinct ordinates; an unattached additional zero of $Y$ costs nothing.

\begin{theorem}[Unknown character lists; \stG]\label{zer:thm:blindlabels}
Suppose $\mathcal Y$ has at most as many entries as $\mathcal X$. If
$\liminf_{T\to\infty}M_{X\to Y}(T)/T=0$, the lists have the same entries with the same multiplicities, up
to permutation. For equal list lengths, every finite set $S$ of primes gives
\begin{equation}\label{zer:eq:witness}
 \liminf_{T\to\infty}\frac{M_{X\to Y}(T)}T\ge\frac1{2^{\abs S+1}\pi}\Bigl(\sum_{p\in S}\abs{S_{\mathcal Y}(p)-S_{\mathcal X}(p)}\frac{\log p}{\sqrt p}-\log\frac{Q_{\mathcal Y}}{Q_{\mathcal X}}\Bigr),
\end{equation}
and for unequal multisets of equal length the right side is strictly positive for some $S$, in either
comparison direction. No conductor bound or list of candidate characters is required.
\end{theorem}
\begin{proof}
If $\mathcal Y$ has fewer characters, $M_{X\to Y}(T)\ge N_X(T)-N_Y(T)\gg T\log T$ by \eqref{zer:eq:coupledcount}.
For equal lengths write $\delta(p)=S_{\mathcal Y}(p)-S_{\mathcal X}(p)$ and $\theta_p=\arg\delta(p)$ when
$\delta(p)\ne0$, arbitrary otherwise, and form the nonnegative trigonometric polynomial
\[
 W(t)=\prod_{p\in S}\bigl(1+\cos(t\log p-\theta_p)\bigr),\qquad0\le W\le2^{\abs S}.
\]
Unique factorization makes its constant coefficient $1$, and its other frequencies are $\log r$ for
$r=\prod_{p\in S}p^{\eps_p}$, $\eps_p\in\set{-1,0,1}$; only $r=p$ and $r=1/p$ can be a prime power or its
reciprocal. For $r=p$ the coefficient is $e^{-i\theta_p}/2$, and the $r=1/p$ sum of a real measure is the
conjugate of the one at $p$. Hence \eqref{zer:eq:coupledlandau} and \eqref{zer:eq:coupledcount} give
\[
 \lim_{T\to\infty}\frac1T\Bigl(\sum_{\gamma\in Y\cap(0,T]}W(\gamma)-\sum_{\gamma\in X\cap(0,T]}W(\gamma)\Bigr)
 =\frac1{2\pi}\Bigl(\log\frac{Q_{\mathcal Y}}{Q_{\mathcal X}}-\sum_{p\in S}\abs{\delta(p)}\frac{\log p}{\sqrt p}\Bigr),
\]
and the expression before the limit is at least $-2^{\abs S}M_{X\to Y}(T)/T$, which gives
\eqref{zer:eq:witness}. Distinct primitive characters induce distinct characters on the unit group of a
common modulus, so unequal multisets give $\delta\ne0$ on some reduced residue class, where $\abs\delta$ is
constant; the reciprocal primes in that class have divergent sum, and $(\log p)/\sqrt p\ge1/p$, so the
prime sum in \eqref{zer:eq:witness} exceeds any prescribed number. Interchanging the lists gives the other
direction.
\end{proof}

For instance $\mathcal X=(\chi_{-4},\chi_5)$ and $\mathcal Y=(\chi_{-3},\chi_5)$ share their even character,
the odd characters take the values $-1$ and $1$ at $7$, and $S=\set7$ gives
\begin{gather*}
 \liminf\frac{M_{X\to Y}(T)}T\ge\frac{2\log7/\sqrt7+\log(4/3)}{4\pi}>0.1399,\\
 \liminf\frac{M_{Y\to X}(T)}T\ge\frac{2\log7/\sqrt7-\log(4/3)}{4\pi}>0.0941,
\end{gather*}
for every pair of Hermitian $2\times2$ couplings; ordinary zero counts have the larger scale $T\log T$.

\begin{corollary}[Sublinear symmetric corruption; \stG]\label{zer:cor:noise}
For any two couplings, with $D(T)=\sum_{0<\gamma\le T}\abs{\nu_A(\gamma)-\nu_B(\gamma)}$,
\[
 \liminf_{T\to\infty}\frac{D(T)}T\ge\sup_{p\text{ prime}}\frac{\log p}{2\pi\sqrt p}\abs{S_{\mathcal X}(p)-S_{\mathcal Y}(p)},
\]
which is positive whenever the label multisets differ, whatever their sizes and conductors. Hence the
positive ordinates after insertions and deletions of total size $o(T)$ through height $T$ determine the
list size, by \eqref{zer:eq:coupledcount}, and the multiset of labels, assuming exact positions of the
retained entries.
\end{corollary}
\begin{proof}
Each changed multiplicity alters a trace sum by at most one in modulus, so sublinear corruption preserves
every limit in \eqref{zer:eq:coupledlandau}; subtract the two formulas at a prime and take the supremum.
If all prime contrasts vanish, use primes in every unit class of a common modulus and orthogonality. The
leading coefficient in \eqref{zer:eq:coupledcount} gives $n$.
\end{proof}

For one-sided contamination the separate prime means can change; the positivity argument of
\cref{zer:thm:blindlabels} treats that case directly.

\paragraph{The coupling at lower order.}
When the label multisets agree, every prime main term agrees, and the coupling appears in a negative
power of $\log t$. For a formal variable $U$, an indeterminate $w\in\C^h$ and the all-ones vector
$\mathbf 1$, both $D_A(U\mathbf 1+w)$ and $D_B(U\mathbf 1+w)$ are monic of degree $n$ in $U$; expand
\begin{equation}\label{zer:eq:couplinglog}
 \log\frac{D_A(U\mathbf 1+w)}{D_B(U\mathbf 1+w)}=\sum_{r\ge1}h_r(w)U^{-r}
\end{equation}
with zero constant term. For distinct labels
$h_r(w)=\tfrac{(-1)^{r+1}}r\tr[(A+\zerdiag w)^r-(B+\zerdiag w)^r]$.

\begin{theorem}[The coupling signal; \stG]\label{zer:thm:couplinggap}
Let $A,B$ be couplings on the same list with $D_A\ne D_B$.
\begin{enumerate}[label=\textup{(\roman*)}]
\item There is a first index $j$ at which $h_j$ is nonconstant, and $j\le n+1$. For distinct labels, if $r$
is the least size of an unequal principal minor, then $j\le\min\set{n+1,2r}$; for Laplacians of graphs with
nonnegative weights $j\le\min\set{n+1,4}$.
\item With $H_A,H_B$ as in \eqref{zer:eq:poissonH} on a fixed line $\Re s=\sigma>1$,
\begin{equation}\label{zer:eq:gapsig}
 H_A(t)-H_B(t)=(\log t)^{-j}f(t)+O((\log t)^{-j-1}),\qquad t\to+\infty,
\end{equation}
where $f$ is a nonzero real uniformly almost periodic function of mean zero.
\item There are an integer $m\ge2$ and $c_m\ne0$ with
\[
 \sum_{0<\gamma\le T}\bigl(\nu_A(\gamma)-\nu_B(\gamma)\bigr)m^{i\gamma}=c_m\frac T{(\log T)^j}+O\Bigl(\frac T{(\log T)^{j+1}}\Bigr).
\]
\end{enumerate}
The constants depend on both matrices, the labels and $m$.
\end{theorem}

\begin{proof}
(i) The series \eqref{zer:eq:couplinglog} converges for large $\abs U$, locally uniformly in $w$, with
termwise derivatives. Its first nonzero index is $k_0=n-\ell_0\le n$, where $\ell_0\le n-1$ is the
$U$-degree of $D_A(U\mathbf 1+w)-D_B(U\mathbf 1+w)$, which is nonzero (evaluate at $U=0$). If $h_{k_0}$ is a
nonzero constant $c$, replacing $w$ by $w+v\mathbf 1$ is the same as replacing $U$ by $U+v$, and the next
power gives $h_{k_0+1}(w+v\mathbf 1)=h_{k_0+1}(w)-k_0cv$; so $j\le n+1$. For distinct labels put
$v_i=(U+w_i)^{-1}$; then $D_A(U\mathbf 1+w)/\prod_i(U+w_i)=1+\sum_{I\ne\varnothing}\det A[I]\prod_{i\in I}v_i$,
and in the logarithm of the ratio all monomials in $v$ of degree below $r$ vanish, the homogeneous degree-$r$
part being $\sum_{\abs I=r}(\det A[I]-\det B[I])\prod_{i\in I}v_i$. At order $U^{-2r}$ a monomial of
$v$-degree $d$ yields a polynomial homogeneous in $w$ of degree $2r-d$, so the coefficient of
$\prod_{i\in I}w_i$ in $h_{2r}$ is $(-1)^r(\det A[I]-\det B[I])$, nonzero for some $r$-set $I$. For
Laplacians the diagonal entries are the one-vertex minors and
$\det A[\set{i,k}]=A_{ii}A_{kk}-\abs{A_{ik}}^2$ with $A_{ik}\le0$ fixes every weight, so $r\le2$.

(ii) Put $D_\ell(s)=iL'(s,\psi_\ell)/L(s,\psi_\ell)$ and $c_\ell=\tfrac i2\log(q_\ell/\pi)$. The squarefree-prime
argument of \cref{zer:lem:algindep} shows that $h_j(c+D(s))$ is a nonconstant Dirichlet series. Write
$U=\tfrac i2\log(s/2)$ and $g_{\psi_\ell}(-i(s-\tfrac12))=U+y_\ell(s)$; Stirling's formula and its derivative
give $y_\ell=c_\ell+D_\ell+O(t^{-1})$ and $y_\ell'=D_\ell'+O(t^{-1})$, $'$ meaning $\dd/\dd s$, with $y,y'$
bounded. The product of the inputs cancels in $P_A/P_B$. Differentiate \eqref{zer:eq:couplinglog} at
$(U,y(s))$: lower coefficients are constant, $U'=O(t^{-1})$, and the differentiated tails converge
uniformly, so
\[
 \frac{P_A'}{P_A}-\frac{P_B'}{P_B}=U^{-j}\frac{\dd}{\dd s}h_j(c+D(s))+O((\log t)^{-j-1})+O(t^{-1}),
\]
and $U=\tfrac i2\log t+O(1)$ gives \eqref{zer:eq:gapsig} with
$f(t)=\Re[(2/i)^j\tfrac{\dd}{\dd s}h_j(c+D(s))]_{s=\sigma+it}$. The series inside has only frequencies
$-\log k$, $k\ge2$, converges absolutely on $\Re s>1$ after differentiation, and has a nonzero coefficient;
taking the real part cannot cancel it against its conjugate, which sits at the opposite frequency.

(iii) Take $\sigma=2$ and fix $m$ with nonzero coefficient $a_m$ in $\tfrac{\dd}{\dd s}h_j(c+D(s))=\sum_ka_kk^{-s}$.
In the real part of $(2/i)^j\sum_ka_kk^{-2-it}$ the coefficient at frequency $-\log m$ is
$(2/i)^ja_mm^{-2}/2$, the conjugate series having only positive frequencies. With $v(t)=(\log t)^{-j}$, each
nonresonant integral $\int_3^Tv(t)e^{i\lambda t}\,\dd t$ is $O_j(1/\abs\lambda)$ by integration by parts, the
nonzero frequency gaps being bounded below for fixed $m$, and the resonant integral is
$T/(\log T)^j+O(T/(\log T)^{j+1})$. Absolute summability gives
\[
 \int_0^Te^{it\log m}(H_A-H_B)(t)\,\dd t=\frac{(2/i)^ja_m}{2m^2}\frac T{(\log T)^j}+O\Bigl(\frac T{(\log T)^{j+1}}\Bigr),
\]
and \eqref{zer:eq:unsmooth} for both couplings, whose error is absorbed, gives (iii) with
$c_m=(2/i)^ja_m/(2\pi\sqrt m)$.
\end{proof}

By the triangle inequality, (iii) gives each one-sided multiplicity discrepancy on $(0,T]$ of order at
least $T/(\log T)^j$, and positivity in \eqref{zer:eq:gapsig} gives the same on windows $[T/2,3T]$;
\cref{zer:thm:localcoupled} below improves both to $\gg T$ distinct ordinates.

\begin{example}[Orders $n+1$ and $4$]\label{zer:ex:couplingorders}
For $n\ge3$ give an $n$-cycle edge signs $+1$ and zero diagonal, and reverse one edge sign for $B$, with the
same labels. Every proper induced subgraph is a forest, on which a diagonal sign change removes any
prescribed edge signs, so all proper principal minors agree; in the full determinant only the two
orientations of the cycle are affected, so $D_A-D_B=C=4(-1)^{n-1}$. The coefficient of $U^{n-1}$ in
$D_B(U\mathbf 1+w)$ is $\sum_iw_{j(i)}$, so \eqref{zer:eq:couplinglog} begins
$CU^{-n}-C(\sum_iw_{j(i)})U^{-(n+1)}$: the constant at order $n$ contributes only through $U'=O(1/t)$, and
$j=n+1$ even with distinct labels, the last possible order. For the triangle with the quadratic characters
of discriminants $-3,-4,5$, whose values at $2$ are $-1,0,-1$, one has $C=4$, $h_4=-4(w_1+w_2+w_3)$, $a_2=8i(\log2)^2$, and
\[
 \sum_{0<\gamma\le T}(\nu_A(\gamma)-\nu_B(\gamma))2^{i\gamma}=\frac{64i(\log2)^2}{\pi\sqrt2}\frac T{(\log T)^4}+O\Bigl(\frac T{(\log T)^5}\Bigr),
\]
while the prime traces and all proper principal minors of the pair agree. The principal-minor identity is
verified in exact arithmetic for $3\le n\le8$ in
\begin{center}\artifact{coupled-dirichlet-inverse-spectrum/checks/cycle_traces.py}\\
(record \artifact{coupled-dirichlet-inverse-spectrum/checks/cycle_result.json}).\end{center}

For order four with positive weights and equal ordinary spectra, attach in cyclic order the characters
of discriminants $-3,-4,5,8$ (two odd, two even) to a $4$-cycle with conductances $1,2,1,2$, and let $B$
interchange $1$ and $2$ on every edge. The Laplacians are conjugate by a cyclic permutation, and
$p_A(y)-p_B(y)=3(y_1-y_3)(y_2-y_4)$ for $p_A(y)=\det(A+\zerdiag y)$: all diagonal entries are $3$, the
complementary quadratic coefficient is $9-w^2$ on an edge and $9$ on a nonedge, the three-vertex minors of
a connected Laplacian agree and both full determinants vanish. The difference is invariant under a
common translation, so \eqref{zer:eq:couplinglog} begins at $U^{-4}$ with the nonconstant coefficient
$3(y_1-y_3)(y_2-y_4)$, and $j=4=2r$. All sixteen coefficients are checked in integer arithmetic in
\artifact{dirichlet-coupled-zeros-and-prime-variances/coupled_zeros/verify_cycle.py}.
\end{example}

\paragraph{Distinct discrepancies in short windows.}
For multisets $X,Y$ of real ordinates and an interval $I$ put
\[
 d_+(I)=\#\set{\gamma\in I:\operatorname{mult}_X(\gamma)>\operatorname{mult}_Y(\gamma)},\qquad
 d_-(I)=\#\set{\gamma\in I:\operatorname{mult}_X(\gamma)<\operatorname{mult}_Y(\gamma)},
\]
so that a shared ordinate contributes one discrepancy if its multiplicities differ.

\begin{theorem}[Distinct discrepancies in short windows; \stG]\label{zer:thm:localcoupled}
Compare $X=Z_{\mathcal X,A}$ and $Y=Z_{\mathcal Y,B}$ for lists of the same size. Suppose either that their
character multisets differ and each list consists of distinct characters, or that they agree and, after matching the labels, $D_A\ne D_B$ (for distinct
labels: some principal minor differs). There are $c_0,C,T_0>0$ depending on the two models such that for
every $T\ge T_0$ and $C\sqrt{\log\log T}\le H\le T/2$,
\[
 \min\set{d_+([T-H,T+H]),d_-([T-H,T+H])}\ge c_0H ,
\]
and the interval even contains $\gg H$ distinct ordinates, in increasing order, on which the direction of
excess alternates. For every $b_0>1/\pi$ each direction occurs at least once in
$[T-b_0\log\log\log T,T+b_0\log\log\log T]$ for all sufficiently large $T$.
\end{theorem}

\begin{proof}
Let $\nu=X-Y$ as a signed point measure, choose $\sigma>1$, $b=\sigma-\tfrac12$, $K_b(t)=b/(b^2+t^2)$, and
write $H_b=K_b*\nu=H_A-H_B$ by \eqref{zer:eq:poissonH}. In the same-label case \eqref{zer:eq:gapsig} gives
$H_b(t)=(\log t)^{-j}f_\sigma(t)+O((\log t)^{-j-1})$ with $f_\sigma$ real uniformly almost periodic, nonzero
and of mean zero; the nonzero Dirichlet coefficient remains nonzero on every such line. For unequal
multisets set $j=0$ and subtract \eqref{zer:eq:cline} on the line $\Re s=\sigma$, which holds there by the
same proof:
\[
 f_\sigma(t)=\tfrac12\log(Q_{\mathcal X}/Q_{\mathcal Y})+\Re\sum_{k\ge2}\Lambda(k)\bigl(S_{\mathcal Y}(k)-S_{\mathcal X}(k)\bigr)k^{-\sigma-it}.
\]
In both cases $H_b(t)=O(\log(2+\abs t))$ globally.

Use $\widehat w(\xi)=\int w(t)e^{-i\xi t}\,\dd t$ and two even unit-mass densities: a Gaussian
$w(t)=(\lambda/\pi)^{1/2}e^{-\lambda t^2}$, and a normalized $e^{-L\cosh(kt)}$ with $k<\pi/(2b)$. For either,
define an integrable real kernel $h$ by $\widehat h(\xi)=\pi^{-1}e^{b\abs\xi}\widehat w(\xi)$. For the
Gaussian, $\widehat w$ and all its derivatives have Gaussian decay. In the second case choose
$b<\beta<\pi/(2k)$ and shift the Fourier contour for $w(t)$, $tw(t)$, $t^2w(t)$ to height $\pm\beta$, where
they decay doubly exponentially since $\cos(k\beta)>0$; the multiplier and its first two derivatives are
then exponentially bounded at both infinities, continuous and piecewise smooth at zero, and two
integrations by parts on the half-lines give $h(v)=O((1+\abs v)^{-2})$. As $\widehat K_b=\pi e^{-b\abs\xi}$,
$h*K_b=w$. The bound $K_b*\abs\nu(u)\ll\log(2+\abs u)$ permits absolute convolution, so $w*\nu=h*H_b$.
Replacing $\log(t-v)$ by $\log t+O(t^{-1/2})$ for $\abs v\le\sqrt t$, the discarded tails being
$O(t^{-1/2}\log(2+t))$,
\begin{equation}\label{zer:eq:locsmoothed}
 (w*\nu)(t)=(\log t)^{-j}G_w(t)+O((\log t)^{-j-1}),\qquad G_w=h*f_\sigma ,
\end{equation}
with $G_w$ real uniformly almost periodic. For common labels its mean is zero; the Gaussian multiplier never
vanishes, so $G_w$ has a nonzero Bohr coefficient, and for the second family $\widehat w(\xi)\to1$ pointwise as
$L\to\infty$, so a large $L$ preserves any one nonzero coefficient.

For different multisets let $\delta(p)=S_{\mathcal Y}(p)-S_{\mathcal X}(p)$. The mean of $G_w$ is
$\log(Q_{\mathcal X}/Q_{\mathcal Y})/(2\pi)$, its coefficient at $-\log p$ is
$\delta(p)(\log p)\widehat w(\log p)/(2\pi\sqrt p)$, and the other nonconstant frequencies are signed
logarithms of prime powers. With the Riesz products of the proof of \cref{zer:thm:blindlabels}, phases
chosen in the two opposite directions, the means of $G_w$ times the two nonnegative trigonometric
polynomials are
\[
 \frac1{2\pi}\Bigl(\log(Q_{\mathcal X}/Q_{\mathcal Y})\pm\sum_{p\in S}\frac{\log p}{\sqrt p}\abs{\delta(p)\widehat w(\log p)}\Bigr),
\]
the coefficients at $\pm\log p$ of a Riesz product having modulus $\tfrac12$ and the two conjugate terms
adding. Choose $S$ to exceed the conductor term by a strict margin when $\widehat w=1$, then $\lambda$ (or
$L$) large enough to preserve the margin; this gives a positive and a negative value of $G_w$. In the
common-label case the same follows from zero mean, since a nonzero nonnegative uniformly almost-periodic
function has positive mean. So there are $a_0>0$ and $D<\infty$ such that every interval of length $D$
contains points with $G_w>2a_0$ and with $G_w<-2a_0$, by translating one point of each sign by relatively
dense approximate periods; in successive blocks of length $D$ select samples of alternating sign, whose
values in \eqref{zer:eq:locsmoothed} have magnitude at least $a_0(\log t)^{-j}$ at large heights.

For the Gaussian and $I=[T-H,T+H]$, at $t$ in the middle half of $I$ the contribution of atoms outside $I$ is
$O(\log(2+T)e^{-c_1H^2})$, by the unit-interval bound summed over distances $m\ge H/2-1$; for
$H\ge C\sqrt{\log\log T}$ with $C$ large this is $o((\log T)^{-j})$. So the Gaussian sum over the atoms in
$I$ alone has $\gg H$ alternating samples. Multiplying it by $e^{\lambda t^2}$ gives a constant times
$E(t)=\sum_\gamma b_\gamma e^{2\lambda\gamma t}$ with $\operatorname{sign}b_\gamma=\operatorname{sign}\nu(\set\gamma)$.
An exponential sum with increasing real exponents has at most as many distinct real zeros as sign changes
of its ordered nonzero coefficients: with no sign change this is immediate; otherwise take $q$ between two
successive exponents at a sign change, and $E\mapsto e^{qt}(e^{-qt}E)'$ multiplies each coefficient by its
exponent minus $q$, removing exactly one sign change, while Rolle's theorem removes at most one zero. This
is an instance of Schoenberg's variation-diminishing property \cite{Schoenberg1948Variation}. Hence the
ordered atoms of $I$ have $\gg H$ sign changes, giving the alternating sequence and the bound in both
directions.

For the shorter window take $\sigma>1$ close to $1$ and $k<\pi/(2b)$ close to its endpoint so that
$kb_0>1$, and use the second density. The positive and negative sample points can be taken within $D$ of
$T$; for either, the mass from ordinates outside the window is
$O(\log(2+T)\exp\set{-c_2\exp(kb_0\log\log\log T)})=o((\log T)^{-j})$, so the signs persist after
truncation, and positivity of $w$ forces an atom of each sign in the window.
\end{proof}

\begin{corollary}[Blind recovery from positive ordinates; \stG]\label{zer:cor:blindall}
Let $\mathcal X$ have $n$ distinct primitive characters, of arbitrary parities and conductors, and let the
competitor $(\mathcal Y,B)$ have at most $n$ distinct ones. If $\liminf_{T\to\infty}d_+((0,T])/T=0$, the two label
multisets coincide and, after matching, $D_A=D_B$; for distinct labels all corresponding principal minors
agree, and for weighted graph Laplacians with nonnegative weights every edge weight is identified. In
particular this holds if the competitor misses $o(T)$ multiplicities of the positive zeros of
$\mathcal P_A$ through height $T$, whatever its additional zeros.
\end{corollary}
\begin{proof}
For fewer competing labels \eqref{zer:eq:coupledcount} gives $\gg T\log T$ missing multiplicities, and the
unit-interval bound allows at most $O(\log(2+T))$ at each ordinate, so $d_+((0,T])\gg T$. At equal size,
if the multisets differ or the grouped polynomials differ after matching, apply
\cref{zer:thm:localcoupled} with centre $2T/3$ and radius $T/3$. Finally the one- and two-vertex minors
give $A_{ii}$ and $\abs{A_{ij}}$, and the off-diagonal signs of a Laplacian are fixed.
\end{proof}

The vertex bound is sharp for one-sided observation: a coupling of larger size can contain the given one
as a block-diagonal summand with further labels, inheriting every zero, because the determinant and the
product of the vertex completions both factor. With symmetric corruption the size is recovered as well
(\cref{zer:cor:noise}). Under \cref{src:thm:dgrh} the corollary contains the finite-discrepancy
statements on the whole line for unknown lists, and gives order $T$ for each one-sided discrepancy in
place of the order $T/(\log T)^{n+1}$ that a fixed-frequency trace yields. The local theorem is the
coupled analogue of the triple-logarithmic windows for products of ordinary Dirichlet functions in
\cref{rec:thm:bothsigns}; here the informative oscillation first appears in a negative power of $\log T$.

Without \cref{src:thm:dgrh}, the same conclusion holds for complex zero data at equal size. This locates
where the passage from complex zeros to signed real ordinates uses the critical-line input: only to infer
poles from missing real zeros.

\begin{theorem}[Labels from the complex divisor; \stU]\label{zer:thm:complexdivisor}
Let $A,B$ be Hermitian couplings on lists $\mathcal A,\mathcal B$ of primitive characters of the same size $n$,
the conductor-one character and repetitions allowed. If
$\sum_{x\in\C}\max\set{\ord_x\mathcal P_A-\ord_x\mathcal P_B,0}<\infty$, then the
character multisets coincide and, after relabelling, $D_A=D_B$.
\end{theorem}

\begin{proof}
Put $C_A(s)=\mathcal P_A((s-\tfrac12)/i)$ and define $C_B$ in the same way. For a list $\mathcal L$ let
$a_{\mathcal L}$ be the sum of its parities and $h_{\mathcal L}$ the number of copies of the conductor-one
label, and put
\[
 c=\tfrac12\log(Q_{\mathcal B}/Q_{\mathcal A}),\qquad\beta=\tfrac12(a_{\mathcal B}-a_{\mathcal A})+2(h_{\mathcal B}-h_{\mathcal A}),\qquad
 G(s)=\frac{\prod_{\chi\in\mathcal B}L(s,\chi)}{\prod_{\chi\in\mathcal A}L(s,\chi)} .
\]
The hypothesis supplies a monic polynomial $W$ cancelling all poles of $C_B/C_A$ with their orders. Let
$w=\deg W$, $e=w+\beta$ and $E=WC_B/C_A$; as in \cref{zer:lem:invrational}, $E$ is entire of order at most
one. Uniformly for $\Re s\ge2$ and $\abs s\to\infty$ there is a fixed $d\ne0$ with
\begin{equation}\label{zer:eq:blindasym}
 E(s)=d\,e^{cs}s^eG(s)\bigl(1+O(1/\log\abs s)\bigr),
\end{equation}
with the principal logarithm: divide each determinant by the product of its completions, so that both
normalized determinants have leading term $(\tfrac i2\log(s/2))^n$ with relative error $O(1/\log\abs s)$;
in the ratio of completion products the equal numbers of gamma factors cancel their common part,
$\Gamma((s+1)/2)/\Gamma(s/2)=(s/2)^{1/2}(1+O(1/\abs s))$, each conductor-one occurrence contributes two
powers through $s(s-1)$, the conductor exponentials give $e^{cs}$, the polynomial gives
$s^w(1+O(1/\abs s))$, and constants go into $d$.

With $\kappa=\abs c$, boundedness of $G$ on $\Re s\ge2$ gives $\abs{E(s)}\ll(1+\abs s)^ee^{c\Re s}$ there.
Realness of $\mathcal P_A,\mathcal P_B$ gives $\abs{C_j(1-\overline s)}=\abs{C_j(s)}$, hence
$\abs{E(s)}\ll(1+\abs s)^ee^{c(1-\Re s)}$ on $\Re s\le-1$. In the strip consider $E(s)/(s+3)^e$, bounded on
both vertical boundaries; for $0<\alpha<\pi/3$ apply the maximum principle on expanding rectangles to it
times $\exp[-\delta\cos(\alpha(s-\tfrac12))]$, whose exponent has real part at least
$\cos(3\alpha/2)\cosh(\alpha\Im s)$ in the strip, and let $\delta\downarrow0$. So
$\abs{E(s)}\le C(1+\abs s)^e\exp(\kappa\abs{\Re s})$ on $\C$.

Choose $T_j\to\infty$ with $p^{-iT_j}\to1$ for every prime $p$ and $e^{icT_j}\to1$, by compactness of finite
tori and a diagonal choice. Absolute convergence gives $G(s+iT_j)\to G(s)$ locally uniformly on $\Re s>1$,
and the entire functions $T_j^{-e}E(s+iT_j)$ are locally bounded, so a subsequence converges to an entire
$H$ with $\abs{H(s)}\le Ce^{\kappa\abs{\Re s}}$ and $H(s)=d_0e^{cs}G(s)$ on $\Re s\ge2$, $d_0=de^{i\pi e/2}\ne0$.
Write $G(s)=\sum_kb_kk^{-s}$ and, for fixed $k$, $\lambda=c-\log k$. On $\Re s=2$ the vertical mean of
$H(s)e^{-\lambda s}$ is $d_0b_k$; by the rectangle with sides on $\Re s=2$ and $\Re s=-u$ and the global
bound, the same mean holds on $\Re s=-u$ for every $u>0$, so $\abs{d_0b_k}\le Ce^{(\kappa+\lambda)u}$. Letting
$u\to\infty$, $G$ is a Dirichlet polynomial supported on $k\le\max(1,Q_{\mathcal B}/Q_{\mathcal A})$. In
particular its prime coefficients $b_p=S_{\mathcal B}(p)-S_{\mathcal A}(p)$ vanish for all large primes,
so the two character sums agree on every unit class of a common modulus and orthogonality identifies the
multisets. After matching, $H=\mathcal P_B-\mathcal P_A\in\mathcal V$ has $\ord_xH\ge\min(\ord_x\mathcal P_A,\ord_x\mathcal P_B)$,
so $H/\mathcal P_A$ has finitely many poles; by \cref{zer:lem:invrational,zer:lem:outside}, $H=0$.
\end{proof}

\begin{corollary}[Finite windows for a bounded catalogue; \stG]\label{zer:cor:catalogue}
Fix a finite catalogue of primitive characters, a maximum list size $N$, a norm bound $M$ and an integer
$K\ge0$. There is $T<\infty$ such that, for all lists from the catalogue with $n_B\le n_A\le N$ and all
couplings with $\norm A,\norm B\le M$, if at most $K$ zero multiplicities of $\mathcal P_A$ in $[-T,T]$ are
missing from $\mathcal P_B$, then the lists coincide as multisets and $D_A=D_B$ after relabelling.
\end{corollary}
\begin{proof}
If no window excluded unequal multisets uniformly, take counterexamples on $[-j,j]$; finitely many lists
and compactness of the matrix balls give a subsequence with fixed distinct multisets and convergent
matrices, both limiting determinants nonzero with real zeros. Around finitely many distinct limit zeros
of the first, take disjoint discs whose closures contain no other zeros of either limit and whose
boundaries avoid them; Rouch\'e's theorem preserves the counts, so the sum of positive parts of the
differences is at most $K$. Exhausting the zeros gives one-sided discrepancy at most $K$ on $\R$,
contrary to \cref{zer:thm:blindlabels}. Once the multisets agree, take the largest threshold of
\cref{zer:thm:window} over the finitely many lists.
\end{proof}

\paragraph{Decoding and stability.}
Let the character list be known, let $K$ bound deleted and $\ell$ bound inserted multiplicities in the
window $[-T_0,T_0]$ of \cref{zer:thm:window} for $(M,K,\ell)$, a spurious extra copy at a genuine zero
counting as one insertion, and let the received multiset be $Y=Z+J$, where $Z$ consists of the true zeros
of $\mathcal P_A$ in the window after at most $K$ deletions and $\abs J\le\ell$. The decoder receives only
$Y$. Put $d=\prod_j(m_j+1)-1$.

\begin{theorem}[Exact decoding by linear equations; \stG]\label{zer:thm:decode}
Let $\norm A\le M$. Seek a monic polynomial $W$ of degree $\ell$ and polynomials $Q_\alpha$ of degree at most
$\ell$ with
\[
 \Bigl(WE_m+\sum_{\alpha\ne m}Q_\alpha E_\alpha\Bigr)^{(r)}(y)=0\qquad(y\in\supp Y,\ 0\le r<\operatorname{mult}_Y(y)).
\]
The system is consistent, and every solution has $Q_\alpha=c_\alpha(A)W$, so the coefficient of $z^\ell$ in
$Q_\alpha$ is $c_\alpha(A)$. No search over corrupted positions is needed.
\end{theorem}
\begin{proof}
$W=z^{\ell-\abs J}\prod_{y\in J}(z-y)$, with multiplicities, and $Q_\alpha=c_\alpha W$ are feasible. For any
solution, the fitted function minus $W\mathcal P_A$ lies in $V_\ell$ and vanishes to the retained orders along
$Z$, so it misses at most $K$ multiplicities of $\mathcal P_A$ in the window; by \cref{zer:thm:window} it
vanishes identically, and \cref{zer:lem:algindep} gives the identities coefficient by coefficient.
\end{proof}

There are $\abs Y$ equations and $\ell+d(\ell+1)$ unknowns; once exact jet values are available, assembly
and Gaussian elimination take polynomially many arithmetic operations, in a model with exact real or
complex arithmetic and zero tests; nothing is asserted about the bit cost of computing transcendental
ordinates. For distinct labels $d=2^n-1$. Spare error budget may leave $W$ nonunique.

\begin{remark}[Certifying a window]\label{zer:rem:windowsearch}
The affine solution set is found by linear algebra; testing $Q_\alpha=c_\alpha W$ on one solution and the
homogeneous identities on a basis of directions certifies a candidate under the corruption model even in a
shorter window, and enlarging the window with complete observations eventually makes the test succeed.
A sufficient $T_0$ in \cref{zer:thm:window} is also computable, for a fixed list, integer budgets and a
rational bound $M$, by a terminating (potentially very slow) search that needs no oracle for real zero
locations. Enclose the norm ball in finitely many rational boxes of Hermitian matrices, and normalize a
coefficient of greatest modulus of $H\in V_\ell$ to one, the others in $[-1,1]$; every parameter pair
$(A,H)$ has finitely many disjoint rational rectangles, with boundaries free of zeros, on which
$\sum_D\max\set{N_D(\mathcal P_A)-N_D(H),0}>K$, and these counts persist in a parameter neighbourhood.
Enumerate rational parameter boxes, rectangles, boundary meshes and precisions; certify by convergent
interval enclosures that each mesh segment has image in an open disc excluding zero, uniformly on the
box, and compute the winding at a rational representative by a polygon whose vertices are enclosed in the
same discs; dovetail until the retained boxes cover the compact charts, and output any integer larger than
the absolute real parts of all rectangle vertices used. Effective enclosures of the completed functions
on compact sets follow from the theta integrals split by the functional equation, or from Hurwitz-zeta
expressions with Euler--Maclaurin bounds, with derivatives by Cauchy's formula. Compactness gives
termination, and reality of the zeros converts rectangle discrepancies into the real-window bound. The
same search on pairs from a finite catalogue gives \cref{zer:cor:catalogue}. No practical analytic bound
on $T_0$ in terms of the conductors is known.
\end{remark}

For stability fix the labelled list, $M$, $K$ and $\ell$, let $T_0$ be as in \cref{zer:thm:window}, put
$T=T_0+1$, and use Euclidean norms of polynomial coefficient vectors in increasing monomial order; the norm
on $V_\ell$ is well defined by \cref{zer:lem:algindep}. For a finite real multiset $D$ of $b$ points with
multiplicity let $p_D(z)=\prod_{x\in D}(z-x)$, and let $I_Df$ be the polynomial of degree less than $b$ with
$f-I_Df$ divisible by $p_D$ among entire functions ($I_\varnothing f=0$).

\begin{lemma}[Stable interpolation; \stU]\label{zer:lem:stableinterp}
Fix $B_0$ and $U$, and let all nodes lie in $[-U,U]$ with $b\le B_0$. If $D,D'$ have the same size and can be
paired within distance $\delta$, then $\norm{I_Df-I_{D'}f}\le C\delta\max_{\abs w=U+2}\abs{f(w)}$. If
$D\subseteq Y$ as multisets with $\abs Y\le B_0$, then $\norm{I_Df}\le C'\norm{I_Yf}$. Both constants depend only
on $B_0,U$.
\end{lemma}
\begin{proof}
On the circle $\abs w=U+2$,
$I_Df(z)=\frac1{2\pi i}\oint f(w)\frac{p_D(w)-p_D(z)}{(w-z)p_D(w)}\,\dd w$, a polynomial of degree less than $b$
whose difference from $f$ is a holomorphic multiple of $p_D$ near the nodes. Each kernel coefficient has
bounded first derivatives in the node variables on the compact box, since $\abs{p_D(w)}\ge2^b$; this
proves the first assertion, collisions included. For the second, $I_Df=I_D(I_Yf)$ by divisibility, and the
integral is bounded on polynomials of degree less than $B_0$ with unit coefficient norm.
\end{proof}

\begin{lemma}[A uniform floor; \stG]\label{zer:lem:stablefloor}
There is $a>0$ such that for all $\norm A\le M$, every multiset $Z$ obtained from the zeros of $\mathcal P_A$
in $[-T,T]$ by at most $K$ deletions, and every $H=\sum_{\alpha\ne m}h_\alpha E_\alpha\in V_\ell$,
$\norm{I_ZH}\ge a\norm{(h_\alpha)_{\alpha\ne m}}$.
\end{lemma}
\begin{proof}
The number of zeros of $\mathcal P_A$ in $[-T,T]$ is uniformly bounded, by local uniform convergence along
convergent subsequences of matrices and Rouch\'e's theorem on a circle enclosing the interval. If the
inequality failed, take $A_j\to A_*$, unit coefficient vectors converging to one defining $H_*$, and $Z_j$
with $I_{Z_j}H_j\to0$, of common cardinality with ordered nodes converging to $Z_*$; the integral formula
gives $I_{Z_*}H_*=0$. Around each zero of $\mathcal P_{A_*}$ in $[-T_0,T_0]$ take a small disc with real
diameter inside $(-T,T)$ and boundary avoiding the zeros of $\mathcal P_{A_*}$ and $Z_*$; for large $j$ the
counts of $\mathcal P_{A_j}$ there equal those of $\mathcal P_{A_*}$, all real, so $Z_*$ contains all but at
most $K$ multiplicities of the zeros of $\mathcal P_{A_*}$ in $[-T_0,T_0]$, and $p_{Z_*}\mid H_*$ contradicts
the choice of $T_0$.
\end{proof}

\begin{theorem}[Lipschitz stability of the decoder; \stG]\label{zer:thm:stabledecoder}
Let $\norm A\le M$ and $0\le\delta\le1$. Erase at most $K$ zeros of $\mathcal P_A$ in $[-T,T]$, displace each
retained one by at most $\delta$, and insert at most $\ell$ arbitrary points in $[-T-1,T+1]$; let $Y$ be the
result, with no information about the matching. There are $R_0,C$ depending only on the list, $M,K,\ell,T$
such that any minimizer of $\norm{I_Y(WE_m+\sum_{\alpha\ne m}Q_\alpha E_\alpha)}$ over monic $W$ of degree
$\ell$ and $\deg Q_\alpha\le\ell$, with joint coefficient norm at most $R_0$, gives $\widehat c_\alpha=[z^\ell]Q_\alpha$
with
\[
 \bigl\|(\widehat c_\alpha-c_\alpha(A))_{\alpha\ne m}\bigr\|\le C\delta .
\]
The minimization is convex in $\ell+d(\ell+1)$ scalar coefficients. Repeated zeros may split into separate
reported locations, and distinct zeros may be reported at one location.
\end{theorem}
\begin{proof}
Choose $R_0$ to contain all vectors $(W_*,(c_\alpha(A)W_*)_\alpha)$ for $\norm A\le M$ and monic $W_*$ of degree
$\ell$ with roots in $[-T-1,T+1]$; for the actual insertions take $W_*=z^{\ell-\abs J}p_J$. With $Z$ the
retained zeros before displacement, $p_{Z\cup J}$ divides $W_*\mathcal P_A$ and $Z\cup J$ pairs with $Y$
within $\delta$, so \cref{zer:lem:stableinterp} gives this feasible vector residual at most $C_1\delta$, and a
minimizer exists and does at least as well. For its fitted function $G$ and the displaced genuine part
$Z'\subseteq Y$, the second interpolation bound gives $\norm{I_{Z'}G}\le C_2\delta$, and the first, with $G$
uniformly bounded on the circle, gives $\norm{I_ZG}\le C_3\delta$. Since $I_Z(W\mathcal P_A)=0$,
$H=G-W\mathcal P_A=\sum(Q_\alpha-c_\alpha(A)W)E_\alpha$ obeys the same bound, and \cref{zer:lem:stablefloor} gives
$\norm{(Q_\alpha-c_\alpha(A)W)_{\alpha\ne m}}\le(C_3/a)\delta$; take the coefficient of $z^\ell$.
\end{proof}

The constants are existential and give no numerical condition number; the contour representation makes the
bound uniform over collisions, whereas unscaled point-evaluation rows can be conditioned very differently.

\begin{remark}[Decoding experiments and unconditional interpolation]\label{zer:rem:decode}
Unconditionally, by \cref{zer:lem:algindep}, the values of $\mathcal P_A/\prod_iF_i$ at $d$ suitable points
$z=-i(s-\tfrac12)$, $s>1$, determine the coefficients $c_\alpha(A)$ by linear algebra. For
$\chi_5,\chi_8,\chi_{12}$ and the path with weights $2,3$, $100$-digit evaluation at $s=2,\dots,8$ recovered
the seven coefficients to about $5.5\cdot10^{-95}$, and an exhaustive check of the $729$ zero-diagonal
symmetric four-vertex matrices with entries in $\set{-1,0,1}$ found $108$ classes, the same for switching
and for principal minors (\artifact{grh_inverse_spectral_2026-11-27/checks.txt}). For the mixed-parity pair of
the quadratic characters modulo $3$ (odd) and $5$ (even) with $A=\begin{psmallmatrix}2&1\\1&3\end{psmallmatrix}$,
$D_A=X_1X_2+3X_1+2X_2+5$, a $45$-digit computation refined nine zeros, deleted the one near
$10.50530281214$ and inserted $4.25$; the system of \cref{zer:thm:decode} returned the location $4.25$ and the
coefficients $3,2,5$ with normalized residual about $2.4\cdot10^{-45}$ (selected roots only; no completeness
or window length is certified). An exact rational polynomial model with one deletion, three insertions and
one excess multiplicity gave an affine system of rank $15$, locator $z(z+1)(z-87)$ and the same coefficients; scripts and records:
\begin{center}\artifact{coupled-dirichlet-inverse-spectrum/checks/dirichlet_locator.py},\\
\artifact{coupled-dirichlet-inverse-spectrum/checks/exact_locator.py},\\
\artifact{coupled-dirichlet-inverse-spectrum/checks/dirichlet_result.json},\\
\artifact{coupled-dirichlet-inverse-spectrum/checks/exact_result.json}.\end{center}
\end{remark}




\subsection{Transport of infinitely many zeros and local Cauchy laws}\label{zer:sec:transport}
\sources{dirichlet_zero_transport_2026-11-26; graph_coupled_zeros_2026-11-26; dirichlet-sharp-window-transport}

The resistance formula \eqref{zer:eq:resistance} concerns finite polynomials, and an unregularized
sum of the displacements of all zeros of the completed functions diverges. The correct objects are
a bounded crossing count and its Poisson transform: the total displacement has a universal growth
law, every coefficient of its smoothed expansion is a finite matrix trace, and for a second
construction in the squared ordinate the leading coefficient is the drop of the Kirchhoff index of
the graph. The same positive phase gives a local Cauchy law for the logarithmic derivative on value
windows that may shrink arbitrarily fast. The spectral-shift mechanism is classical
\cite{BehrndtGesztesyNakamura2018Spectral,Reiffenstein2022Interlacing}; weak Cauchy limits for the
logarithmic derivative of $\xi$ are due to Sodin and to Najnudel and Nikeghbali
\cite{Sodin2018Critical,NajnudelNikeghbali2026Cauchy}, and Cauchy laws for Herglotz boundary values
to Aizenman and Warzel \cite{AizenmanWarzel2015Cauchy}. On every interval in $[T,2T]$ the displacement
equals an explicit integral up to an error uniform in the dimension, the conductors and both matrices,
which gives a relative Cauchy law for every Borel set of values, a finite part for the infinite transport,
and an expansion to every order on sharp windows. No simplicity, spacing, independence or
pair-correlation hypothesis is used.

\paragraph{Setting.}
Let $F_1,\dots,F_n$ be functions \eqref{zer:eq:phasereal} attached to primitive characters of
conductors $q_i$ (repetitions allowed; $q_i=1$ for $F_\zeta$), and $g_i=F_i'/F_i$. Under
\cref{src:thm:dgrh}, \eqref{zer:eq:logderim} gives $\Im g_i<0$ in the upper half-plane. For a
Hermitian $A$ let $\mathcal P_A$ and $M_A$ be as in \eqref{zer:eq:PA}; $M_A$ is invertible in the upper
half-plane and $\mathcal P_A$ has only real zeros. Fix a Hermitian $B\succeq0$ of rank $r\ge1$ and put,
at real $x$ avoiding the discrete set of input and output zeros,
\begin{equation}\label{zer:eq:shift}
 \sigma_{A,B}(x)=n_+(M_{A+B}(x))-n_+(M_A(x)),\qquad 0\le\sigma_{A,B}\le r .
\end{equation}

\begin{theorem}[Displacement measure and trace identity; \stG]\label{zer:thm:shift}
There is a holomorphic logarithm $\mathcal L_{A,B}$ of $\mathcal P_{A+B}/\mathcal P_A$ on the upper
half-plane with $0<\Im\mathcal L_{A,B}<r\pi$ and $\mathcal L_{A,B}(iy)\to0$ as $y\to\infty$. Its boundary
imaginary part is $\pi\sigma_{A,B}$ almost everywhere,
\[
 \Im\mathcal L_{A,B}(x+iy)=\int_\R\frac{y\,\sigma_{A,B}(t)}{(t-x)^2+y^2}\,\dd t,
\]
and $\dd\sigma_{A,B}=\dd\nu_A-\dd\nu_{A+B}$ as distributions, where $\nu_A$ is the zero-counting measure
of $\mathcal P_A$ with multiplicity. Hence for every $\varphi\in C^1_c(\R)$
\[
 \sum_{\mathcal P_{A+B}(\beta)=0}\ord_\beta(\mathcal P_{A+B})\varphi(\beta)-\sum_{\mathcal P_A(\alpha)=0}\ord_\alpha(\mathcal P_A)\varphi(\alpha)
 =\int_\R\varphi'(t)\sigma_{A,B}(t)\,\dd t,
\]
common zeros cancelling with their exact multiplicities.
\end{theorem}
\begin{proof}
Write $B=VV^*$ with $V$ of full column rank $r$ and $H(z)=V^*M_A(z)^{-1}V$. Since
$\Im(M^{-1})=-M^{-*}(\Im M)M^{-1}$, $\Im H\succ0$ in the upper half-plane, so the eigenvalues of
$I_r+H$ lie there and $\mathcal L_{A,B}=\tr\log(I_r+H)$, with the principal logarithm, is holomorphic
with imaginary part in $(0,r\pi)$. The determinant lemma gives
$e^{\mathcal L_{A,B}}=\det M_{A+B}/\det M_A=\mathcal P_{A+B}/\mathcal P_A$, and \cref{zer:lem:radial} below gives
$H(iy)=O(1/\log y)$. At a real regular point $H(x)$ is Hermitian, and approaching from above,
$\pi^{-1}\Im\mathcal L_{A,B}(x+i0)=n_-(I_r+H(x))$. The two Schur complements of
$\begin{pmatrix}M_A(x)&V\\V^*&-I_r\end{pmatrix}$ give its inertia as $n_+(M_A)+n_-(I_r+H)$ and as
$n_+(M_A+B)$, which identifies the boundary value with $\sigma_{A,B}$. A bounded harmonic function is
the Poisson extension of its boundary values, which gives the Poisson formula. Near a real $x_0$,
$\mathcal P_{A+B}/\mathcal P_A=(z-x_0)^ka(z)$ with $a$ real analytic and nonzero and
$k=\ord_{x_0}\mathcal P_{A+B}-\ord_{x_0}\mathcal P_A$; the boundary argument of $\log(z-x_0)$ drops by
$\pi$ across $x_0$, so $\sigma_{A,B}$ jumps by $-k$, and integration by parts gives the trace identity.
\end{proof}

For a rank-one $B$ the set where the boundary value of $1+H$ is negative is a disjoint union of
intervals from an old zero to the following new zero, after cancelling common factors, and
$\sigma_{A,B}$ is its indicator; so its integral over a window is the total length of the moving-zero
intervals clipped to the window. A general $B$ is a sum of $r$ rank-one terms whose counts add and
telescope.

\begin{lemma}[Radial logarithmic derivative; \stU]\label{zer:lem:radial}
Put $\kappa_i=\tfrac12\log(q_i/(2\pi))$. For every compact $K$ in the upper half-plane,
uniformly for $z\in K$,
\[
 g_i(Tz)=-i\Bigl(\tfrac12\log T+\kappa_i+\tfrac12\log(-iz)\Bigr)+O_K(T^{-1}),
\]
with the principal logarithm. In particular $g_i(iy)=-i(\tfrac12\log y+\kappa_i)+O(y^{-1})$.
\end{lemma}
\begin{proof}
The functional equation gives $g_i(Tz)=-i\,\Lambda_{\overline\chi}'(s)/\Lambda_{\overline\chi}(s)$ with
$s=\tfrac12-iTz$, which stays in a closed subsector of the right half-plane with $\Re s\to\infty$
linearly in $T$. There $\Lambda'/\Lambda=\tfrac12\log(q_i/\pi)+\tfrac12\psi((s+a_i)/2)+L'/L(s)$, the last
term is exponentially small, and $\psi(w)=\log w+O(\abs w^{-1})$. For $\zeta$ the factor $s(s-1)$ adds
$1/s+1/(s-1)=O_K(T^{-1})$.
\end{proof}

Let
$\mathfrak D_{A,B}(y)=\int_\R y^2\sigma_{A,B}(t)(t^2+y^2)^{-1}\,\dd t=y\Im\mathcal L_{A,B}(iy)$ be the
Poisson-smoothed displacement, $\ell_y=\tfrac12\log y$ and $\mathsf C=\zerdiag(\kappa_i)$.

\begin{theorem}[Complete smoothed expansion; \stG]\label{zer:thm:transportexp}
For every fixed $m\ge1$,
\[
 \frac{\mathfrak D_{A,B}(y)}y=\sum_{k=1}^m\frac{(-1)^{k+1}}{k\ell_y^k}\Im\tr\bigl((\mathsf C+i(A+B))^k-(\mathsf C+iA)^k\bigr)+O(\ell_y^{-m-1}).
\]
In particular
\[
 \frac{\mathfrak D_{A,B}(y)}y=\frac{\tr B}{\ell_y}-\frac{\tr(\mathsf CB)}{\ell_y^2}
 +\frac{\tr(\mathsf C^2B)-\tr(A^2B)-\tr(AB^2)-\frac13\tr(B^3)}{\ell_y^3}+O(\ell_y^{-4}),
\]
the matrices need not commute, the first correction records the conductors, and the coupling $A$ first
appears at the third order. For $A=L_G$ and $B=hvv^*$, $v=e_a-e_b$, the third coefficient is
$h(\kappa_a^2+\kappa_b^2)-h\norm{L_Gv}^2-2h^2v^*L_Gv-\tfrac83h^3$.
\end{theorem}
\begin{proof}
By \cref{zer:lem:radial}, $M_A(iy)=-i\ell_y\bigl(I+(\mathsf C+iA)/\ell_y+O((y\ell_y)^{-1})\bigr)$. Use the matrix
logarithm series at $I$, take traces and subtract the expressions for $A+B$ and $A$; the factors
$(-i\ell_y)^n$ cancel, both branches tend to zero as the normalization of $\mathcal L_{A,B}$ requires, and
the arithmetic remainder costs $O((y\ell_y)^{-1})$. For the first terms,
$\Im\tr(\mathsf C+iA)=\tr A$, $\Im\tr(\mathsf C+iA)^2=2\tr(\mathsf CA)$ and
$\Im\tr(\mathsf C+iA)^3=3\tr(\mathsf C^2A)-\tr(A^3)$, by cyclicity.
\end{proof}

With the zeta input at two vertices joined by conductance $w$, $\mathcal P_w=F'(F'+2wF)$, and with
$a(y)=\xi'(\tfrac12+y)/\xi(\tfrac12+y)>0$ the identity $g(iy)=-ia(y)$ reduces the trace formula to
\[
 \mathfrak D_{w,h}(y)=y\Bigl[\arctan\frac{2(w+h)}{a(y)}-\arctan\frac{2w}{a(y)}\Bigr],
\]
which can be evaluated for $y>\tfrac12$ from
$a(y)=1/s+1/(s-1)-\tfrac12\log\pi+\tfrac12\psi(s/2)+\zeta'(s)/\zeta(s)$, $s=y+\tfrac12$, without locating
any zero. The floating-point diagnostics of the normalized transforms, and exact checks of the trace
identity on polynomial inputs with repeated and shared roots, are in
\artifact{dirichlet_zero_transport_2026-11-26/code/verify_transport.py} (record
\artifact{dirichlet_zero_transport_2026-11-26/data/verification.json}). The expansion is for the
smoothed displacement; its sharp-window counterpart, with a real shift of $A$, is
\cref{zer:cor:sharpexp}, and no microscopic spacing law follows.

\paragraph{The squared ordinate.}
Suppose the inputs are even, write $\ord_0F_i=2m_i$ and $F_i(z)=z^{2m_i}G_i(z^2)$, so that $G_i(0)\ne0$;
this holds for $\xi$ and for real primitive characters. Let $G_0,G_1$ be connected weighted graphs on
$n\ge2$ vertices with $L_{G_1}-L_{G_0}\succeq0$ (for instance after adding edges or increasing
conductances), with spanning-tree sums $\tau_j$, effective resistances $R^{(j)}_{ab}$ and Kirchhoff
indices $\mathrm{Kf}_j=\sum_{a<b}R^{(j)}_{ab}=n\tr L_{G_j}^+$; Rayleigh monotonicity gives
$\Delta\mathrm{Kf}=\mathrm{Kf}_0-\mathrm{Kf}_1\ge0$, with strict inequality if $L_{G_1}\ne L_{G_0}$.
Put $h_i=G_i'/G_i$, $Q_j(w)=\det(\zerdiag(G_i(w))L_{G_j}+\zerdiag(G_i'(w)))$ and
$\eta(x)=n_+(L_{G_1}+\zerdiag h_i(x))-n_+(L_{G_0}+\zerdiag h_i(x))$.

\begin{theorem}[Infinite electrical transport; \stG]\label{zer:thm:squared}
The functions $Q_j$ have only real zeros, bounded below, and $\eta$ is supported on a half-line. If
$L_{G_1}\ne L_{G_0}$, then
\[
 \int_{-\infty}^X\eta(u)\,\dd u\sim\frac{\Delta\mathrm{Kf}}{4\pi n}\sqrt X\log X,
\]
for every sufficiently large $R$ one has the convergent identity, with an unambiguous positive real
logarithm,
\[
 \log\frac{Q_1(-R)/\tau_1}{Q_0(-R)/\tau_0}=\int_\R\frac{\eta(u)}{R+u}\,\dd u,
\]
and for fixed $u_0>0$ the displacement measured back in positive ordinary ordinates satisfies
$\int_{u_0}^{T^2}\eta(u)(2\sqrt u)^{-1}\,\dd u\sim\Delta\mathrm{Kf}\,(\log T)^2/(8\pi n)$. For one added edge,
$L_{G_1}=L_{G_0}+hbb^{\mathsf T}$ with $b=e_a-e_b$, one has
$\Delta\mathrm{Kf}=nh\norm{L_{G_0}^+b}^2/(1+hR_e)$ with $R_e=b^{\mathsf T}L_{G_0}^+b$.
\end{theorem}
\begin{proof}
The zeros of $G_i$ are the $\gamma^2$, $\gamma>0$ a zero of $F_i$; $G_i$ has order at most $\tfrac12$ and the
genus-zero product $G_i(w)=G_i(0)\prod_{\gamma>0}(1-w/\gamma^2)$, so $h_i(w)=\sum_{\gamma>0}(w-\gamma^2)^{-1}$
and the half-plane argument of \cref{zer:thm:shift} applies. At $w=-R$,
\[
 h_i(-R)=\frac{g_i(i\sqrt R)}{2i\sqrt R}+\frac{m_i}R=-\frac{\log R}{8\sqrt R}-\frac{\kappa_i}{2\sqrt R}+O(R^{-1})
\]
by \cref{zer:lem:radial}. With $a_R=\log R/(8\sqrt R)$, for large $R$ the matrix
$L_{G_j}+\zerdiag h_i(-R)$ is positive definite on $\mathbf 1^\perp$ and has a negative Schur complement
in the constant direction, so it has exactly one negative eigenvalue; the same holds below a fixed
$u_*$, so $Q_j$ has no zeros and $\eta=0$ there.

For variables $x_i$, the matrix-tree and two-root forest identities give
$\det(L_{G_j}+\zerdiag x)/\tau_j=\sum_ix_i+\sum_{a<b}R^{(j)}_{ab}x_ax_b+O(\norm x^3)$, the coefficient of
$x_ax_b$ being $\det L_{G_j}[[n]\setminus\set{a,b}]$, which is $\tau_jR^{(j)}_{ab}$ by the cofactor formula.
Substituting $x_i=h_i(-R)$, the linear terms are $-na_R(1+O(1/\log R))$, the quadratic terms differ by
$-\Delta\mathrm{Kf}\,a_R^2(1+O(1/\log R))$, and $\prod_iG_i(-R)$ cancels, so
\[
 f(-R):=\log\frac{Q_1(-R)/\tau_1}{Q_0(-R)/\tau_0}=\frac{\Delta\mathrm{Kf}}na_R+O\Bigl(\frac{a_R}{\log R}+a_R^2\Bigr)\sim\frac{\Delta\mathrm{Kf}}{8n}\frac{\log R}{\sqrt R}.
\]
As in \cref{zer:thm:shift}, $f(z)=\tr\log(I+V^*(L_{G_0}+\zerdiag h_i(z))^{-1}V)-\log(\tau_1/\tau_0)$ has bounded
nonnegative imaginary part with boundary value $\pi\eta$, so $f'(z)=\int\eta(u)(u-z)^{-2}\,\dd u$.
Integrating over $[-S,-R]$, below the support of $\eta$, and letting $S\to\infty$ with $f(-S)\to0$ gives
the Stieltjes identity. After translating the support by $u_*$, the Stieltjes form of Karamata's
theorem \cite[Theorem~A.7 and Remark~A.8]{LangerWoracek2025Karamata} with exponent $\tfrac12$ and slowly
varying factor $\log R$ gives the first asymptotic: if $\nu\ge0$ on $[0,\infty)$ has
$\int(R+u)^{-1}\,\dd\nu\sim cR^{\alpha-1}L(R)$, $0<\alpha<1$, then
$\nu([0,X])\sim c\sin(\pi\alpha)X^\alpha L(X)/(\pi\alpha)$. For this converse, the rescaled measures
$\nu(RE)/(RS(R))$ have transforms tending to $s^{\alpha-1}$ with uniformly small tails, every
subsequential vague limit has Stieltjes transform $s^{\alpha-1}$, and uniqueness identifies it as
$\sin(\pi\alpha)u^{\alpha-1}\dd u/\pi$. Finally, with $U(X)=\int_{-\infty}^X\eta$,
$\int_{u_0}^X\dd U/(2\sqrt u)=U(X)/(2\sqrt X)+\tfrac14\int_{u_0}^XU(u)u^{-3/2}\,\dd u+O(1)$, which gives the
last asymptotic with $X=T^2$. The edge formula follows from the rank-one inverse identity on
$\mathbf 1^\perp$.
\end{proof}

For the path on three vertices with unit conductances completed to a triangle, $\tau_0=1$, $\tau_1=3$,
$\mathrm{Kf}_0=4$, $\mathrm{Kf}_1=2$; with $\xi$ at every vertex, $U(X)\sim\sqrt X\log X/(6\pi)$, and with
$m=h_i(-R)$ the normalized ratio is exactly $(m+3)/(3(m+1))$. The ordinary and squared constructions
define different entire functions: replacing $F_i'$ by $G_i'$ changes the coupling.

\paragraph{Sharp windows.}
The displacement has a sharp-interval form with a dimension-free endpoint error, from which the
macroscopic transport law and a local Cauchy law for the logarithmic derivative follow as
corollaries. Keep the setting above with $B\succeq0$, now allowed to be zero, and put, for $t>2\pi/q$,
\[
 d_q(t)=\tfrac12\log\frac{qt}{2\pi},\qquad p_{q,t}(v)=\frac1\pi\frac{d_q(t)}{(v+\pi/4)^2+d_q(t)^2},\qquad
 p^*_{q,T}(v)=\frac1\pi\frac{d_q(T)}{v^2+d_q(T)^2},
\]
two probability densities; write $p(E)=\int_Ep$. For a diagonal $Z$ with $\Im Z\prec0$ define
\begin{equation}\label{zer:eq:phaseTheta}
 \Theta(Z)=\int_0^1\tr\bigl(B(A+uB+Z)^{-1}\bigr)\,\dd u,
\end{equation}
a logarithm of $\det(A+B+Z)/\det(A+Z)$ fixed for all ranks; $\Theta(\zerdiag g_i(z))=\mathcal L_{A,B}(z)$ in the
upper half-plane by \cref{zer:thm:shift}. Let $D(t)=\zerdiag(d_{q_i}(t))$, $D=D(T)$, $d=\min_iD_{ii}$, and
\[
 \Omega=\Im\Theta(-iD),\qquad \rho(t)=\pi^{-1}\Im\Theta\bigl(-\tfrac\pi4I-iD(t)\bigr),
\]
so that $\Omega>0$ if $B\ne0$. With $A'=D^{-1/2}AD^{-1/2}$, $B'=D^{-1/2}BD^{-1/2}$ and $X=A'+uB'$, congruence
invariance of \eqref{zer:eq:phaseTheta} gives
\begin{equation}\label{zer:eq:Omega}
 \Omega=\int_0^1\tr\bigl(B'(I+X^2)^{-1}\bigr)\,\dd u=\tr\bigl(\arctan(A'+B')-\arctan A'\bigr).
\end{equation}

\begin{theorem}[Sharp interval estimate; \stG]\label{zer:thm:sharpinterval}
There are absolute constants $T_0,C$ such that for all $T\ge T_0$ and all intervals $I\subseteq[T,2T]$,
\[
 \Bigl|\int_I\sigma_{A,B}(t)\,\dd t-\int_I\rho(t)\,\dd t\Bigr|\le C\frac\Omega{\log d},\qquad
 \pi\rho(t)=\Omega\bigl(1+O(d^{-1})\bigr)\quad(T\le t\le2T),
\]
the implied constant also absolute: nothing depends on $n$, the characters, the conductors, $A$ or $B$,
which may all vary with $T$.
\end{theorem}

For fixed data and $B\ne0$, $\Omega\sim2\tr(B)/\log T$, so the error is $O(1/(\log T\log\log T))$ for every
endpoint pair in a dyadic interval; uniformity in $A$ and $B$ is what gives the much finer error for a rare
scalar event below. The proof uses three estimates. Write $Z'=D^{-1/2}ZD^{-1/2}$.

\begin{lemma}[Comparing resolvents; \stU]\label{zer:lem:resolvent}
If $\norm{Z_k'+iI}\le\tfrac12$ for $k=1,2$, then $\abs{\Theta(Z_1)-\Theta(Z_2)}\le4\Omega\norm{Z_1'-Z_2'}$; in
particular $\Im\Theta(Z_1)\ll\Omega$.
\end{lemma}
\begin{proof}
Let $R_*=(X-iI)^{-1}$ and $R_k=(X+Z_k')^{-1}$. Factoring the inverse with $\norm{R_*}\le1$ gives
$\norm{R_kB'^{1/2}}_{\mathrm{HS}}\le2\norm{R_*B'^{1/2}}_{\mathrm{HS}}$, and the adjoints obey the same bound since
$R_*$ is normal with $R_*^*R_*=(I+X^2)^{-1}$. In $R_1-R_2=R_1(Z_2'-Z_1')R_2$ apply the Hilbert--Schmidt
Cauchy--Schwarz inequality with $B'^{1/2}$ on the two outside factors; the square of the reference norm
integrates to $\Omega$ by \eqref{zer:eq:Omega}. For the last assertion take $Z_2=-iD$.
\end{proof}

Put $G(z)=\zerdiag(g_i(z))$ and, with the principal logarithm $\zerLog$,
\[G_0(z)=-i\,\zerdiag\bigl(\tfrac12\log(q_i/(2\pi))+\tfrac12\zerLog(-iz)\bigr),\]
whose boundary value at $t>0$ is
$-\tfrac\pi4I-iD(t)$, the real shift coming from the argument of $-it$. Let $Y=2+4\log(2T)$.

\begin{lemma}[A rectangle estimate near the real axis; \stG]\label{zer:lem:strip}
There are absolute $c,C_0>0$ such that for $T\ge T_0$, $T\le t\le2T$, $h=C_0/\log d$ and $z=t+iy$,
\[
 \norm{D^{-1/2}(G-G_0)(z)D^{-1/2}}\ll\begin{cases}e^{-cy\log d}+d^{-1}&(h\le y\le1),\\(T^{-1}+2^{-y})/d&(1\le y\le Y).\end{cases}
\]
Each normalized matrix $G(z)$, $G_0(z)$ is within $\tfrac12$ of $-iI$ in norm for $h\le y\le Y$, and so is
$G_0$ down to $y=0$.
\end{lemma}
\begin{proof}
Set $s=\tfrac12-iz$. The functional equation gives $g_j(z)=-i\Lambda'_{\overline\chi_j}(s)/\Lambda_{\overline\chi_j}(s)$, with
$\xi$ in the zeta case. For $y\ge1$ the logarithmic derivative of the Euler product is $O(2^{-y})$ uniformly
in the character, since $\sum_{k\ge2}\Lambda(k)k^{-3/2}<\infty$, the gamma part is
$\tfrac12\log(q_j/\pi)+\tfrac12\psi((s+a_j)/2)$, and $\psi(w)=\zerLog w+O(1/\abs w)$; so
$g_j=(G_0)_{jj}+O(\abs z^{-1}+2^{-y})$, the rational terms of $\xi$ obeying the same bound. This is the second
case, and on the rectangle $G_0+iD=O(1)$ uniformly since $Y=O(\log T)$.

With $d_j=D_{jj}$, the Cayley transform $w_j=(g_j+id_j)/(g_j-id_j)$ is holomorphic with modulus at most one above
the real axis, its real poles being removable, and $\abs{w_j(x+i)}\le K/d_j$ for $\abs{x-t}\le1$ by the
second case. On the rectangle $\abs{x-t}\le1$, $0\le y\le1$ the harmonic function
$v=\cos(\pi(x-t)/2)\sinh(\pi y/2)/\sinh(\pi/2)$ vanishes on the sides and bottom and lies in $[0,1]$ on the top,
so the maximum principle applied to $\log\abs{w_j}+v\log(d_j/K)$ gives $\abs{w_j(t+iy)}\le\exp(-cy\log d_j)$. With
$C_0$ large, inverting the Cayley transform gives $\abs{g_j/d_j+i}\le2\abs{w_j}/(1-\abs{w_j})\le\tfrac12$ for
$h\le y\le1$, and the first case. Increasing $T_0$ so that $2h<1$ gives the remaining norm assertions.
\end{proof}

The smoothing is removed by the kernel
\[
 K_h(s)=2P_h(s)-P_{2h}(s)=\frac{6h^3}{\pi(s^2+h^2)(s^2+4h^2)},\qquad P_y(s)=\frac y{\pi(s^2+y^2)},
\]
which is nonnegative with integral one and tail $\int_r^\infty K_h\ll\min(1,(h/r)^3)\ll hP_h(r)$. For
$I=[a,b]$, a change between $\mathbf 1_I(s)$ and $\mathbf 1_I(s-u)$ requires crossing $a$ or $b$, so for every
bounded nonnegative $f$
\begin{equation}\label{zer:eq:clip}
 \Bigl|\int_I(f-K_h*f)\Bigr|\ll h\bigl((P_h*f)(a)+(P_h*f)(b)\bigr);
\end{equation}
the extra decay of $K_h$ is essential in this comparison with $P_h$.

\begin{proof}[Proof of \cref{zer:thm:sharpinterval}]
Take $h$ from \cref{zer:lem:strip}. By the Poisson formula of \cref{zer:thm:shift},
$(P_y*\sigma_{A,B})(x)=\pi^{-1}\Im\Theta(G(x+iy))$, so \eqref{zer:eq:clip} and \cref{zer:lem:resolvent} show that
replacing $\sigma_{A,B}$ by $K_h*\sigma_{A,B}$ costs $O(h\Omega)$; its integral is $\pi^{-1}$ times the imaginary
part of the horizontal integrals of $\Theta(G)$ at heights $h,2h$ with coefficients $2,-1$. For
$J(z)=\Theta(G(z))-\Theta(G_0(z))$ and $\eta\in\set{h,2h}$, integrate $J$, holomorphic there, around
$[a,b]+i[\eta,Y]$. By \cref{zer:lem:resolvent,zer:lem:strip} each vertical side costs
\[
 C\Omega\Bigl(\int_h^1(e^{-cy\log d}+d^{-1})\,\dd y+d^{-1}\int_1^Y(T^{-1}+2^{-y})\,\dd y\Bigr)\ll\Omega/\log d,
\]
and the top costs $O(\Omega/d)$ since $b-a\le T$ and $T2^{-Y}\le1$; Cauchy's theorem bounds the bottom
integral by $O(\Omega/\log d)$. This bounds an integral of the arithmetic remainder over a long interval by
its values near the endpoints. Finally $\norm{D^{-1/2}(G_0(t+i\eta)-G_0(t))D^{-1/2}}\ll\eta/(Td)$, so the
integrated change is $O(\Omega\eta/d)$; the coefficients $2,-1$ sum to one, leaving $\int_I\rho$. The pointwise
formula follows from \cref{zer:lem:resolvent}, since $\norm{D^{-1/2}(G_0(t)+iD)D^{-1/2}}\ll1/d$.
\end{proof}

\begin{corollary}[Universal macroscopic transport; \stG]\label{zer:thm:transport}
As $T\to\infty$, the measures $(\log T)\sigma_{A,B}(Tu)\,\dd u$ converge vaguely to
$(2\tr B/\pi)\,\dd u$. Consequently, for fixed
$a<b$,
\[
 \int_{aT}^{bT}\sigma_{A,B}(x)\,\dd x\sim\frac{2\tr B}\pi(b-a)\frac T{\log T},\qquad
 \int_{-T}^T\sigma_{A,B}(x)\,\dd x\sim\frac{4\tr B}\pi\,\frac T{\log T}.
\]
For an edge of added conductance $h$, $B=h(e_a-e_b)(e_a-e_b)^*$ and the last coefficient is $8h/\pi$.
The leading coefficient depends neither on $A$ nor on the characters, conductors or spacings, and
the unregularized total displacement is infinite.
\end{corollary}
\begin{proof}
Fix the data; $B\ne0$ since $r\ge1$. In \eqref{zer:eq:Omega}, $A'$ and $B'$ are $O(1/\log T)$, so
$\Omega=\tr B'(1+O((\log T)^{-2}))\sim2\tr B/\log T$, and the pointwise part of \cref{zer:thm:sharpinterval}
gives $\rho(t)=(2\tr B/\pi)(1+o(1))/\log t$ as $t\to\infty$. Fix $t_0\ge T_0$ and cover $[t_0,X]$ by the
intervals $[2^kt_0,2^{k+1}t_0]$, the last one truncated. On each, \cref{zer:thm:sharpinterval} replaces
$\sigma_{A,B}$ by $\rho$ at a cost $O(\Omega/\log d)=O(1)$, and there are $O(\log X)$ of them, so
\[
 \int_{t_0}^X\sigma_{A,B}=\int_{t_0}^X\rho+O(\log X)=\frac{2\tr B}\pi\,\frac X{\log X}\bigl(1+o(1)\bigr)\qquad(X\to\infty).
\]
For negative heights use the conjugate list $\widetilde F_i(z)=F_i(-z)$, which consists of functions
\eqref{zer:eq:phasereal} of the conjugate characters, with $F_\zeta$ unchanged, the same conductors and
logarithmic derivatives $-g_i(-z)$. At a real $x$ avoiding the discrete exceptional set,
$\widetilde M_{A'}(x)=-M_{-A'}(-x)$ is invertible for every Hermitian $A'$ in question, so
$n_+(\widetilde M_{A'}(x))=n-n_+(M_{-A'}(-x))$, and with $A'=-A-B$ this gives
$\sigma_{A,B}(-x)=\widetilde\sigma_{-A-B,B}(x)$, the displacement of the conjugate list for the Hermitian
matrices $-A-B$ and $B\succeq0$. The previous step applies to these data, with the same $B$, so
$\int_{-X}^{-t_0}\sigma_{A,B}$ has the same asymptotic. Since $0\le\sigma_{A,B}\le r$, the range $\abs x\le t_0$
contributes $O(t_0)$. For fixed $a<b$ the integral over $[aT,bT]$ is a sum or difference of these
one-sided integrals at $X=\abs aT$ and $X=\abs bT$, up to $O(t_0)$, which gives both displays. So the
masses of $[a,b]$ under the measures $(\log T)\sigma_{A,B}(Tu)\,\dd u$ tend to $(2\tr B/\pi)(b-a)$; they are
bounded on each compact set, and approximating continuous compactly supported functions uniformly by
step functions gives the vague convergence. For an edge $\tr B=h\norm{e_a-e_b}^2=2h$. Nothing in the
limit depends on $A$ or on the inputs, and $\int_\R\sigma_{A,B}=\infty$.
\end{proof}

\begin{theorem}[Uniform local Cauchy law; \stG]\label{zer:thm:uniformcauchy}
Let $F$ be one of the functions \eqref{zer:eq:phasereal}, of conductor $q$, $g=F'/F$, and $d=d_q(T)$. For all
sufficiently large $T$, every interval $I\subseteq[T,2T]$ of length $H$ and every Borel set $E\subseteq\R$,
\[
 \Bigl|\,\abs{\set{t\in I:g(t)\in E}}-\int_Ip_{q,t}(E)\,\dd t\Bigr|\le\frac C{\log d}\,p^*_{q,T}(E),
\]
with absolute $C$ and threshold. In particular, if $H>0$ and $p^*_{q,T}(E)>0$, the probability that $g(t)\in E$
for uniform $t\in I$ is $p^*_{q,T}(E)\bigl(1+O(d^{-1}+(H\log d)^{-1})\bigr)$.
\end{theorem}
\begin{proof}
For $E=(u,v)$ take $n=1$, $A=-v$, $B=v-u$; outside a discrete set $\sigma_{A,B}=\mathbf 1_{\set{u<g<v}}$, and direct
integration of \eqref{zer:eq:phaseTheta} gives $\Omega=\pi p^*_{q,T}((u,v))$ and $\rho(t)=p_{q,t}((u,v))$. Apply
\cref{zer:thm:sharpinterval}. The occupation and model measures are finite and give interval endpoints
zero mass, $g$ being nonconstant meromorphic; by additivity the two inequalities extend to finite disjoint
unions of intervals and then by the monotone class theorem to Borel sets. The pointwise part of
\cref{zer:thm:sharpinterval} gives the relative form on intervals, extended in the same way.
\end{proof}

For example, for $F_\zeta$, $I=[T,T+H]$ with $H=(\log\log T)^{-1/2}$ and $E=[e^T,e^T+e^{-T}]$, the set of
$t\in I$ with $g(t)\in E$ has length asymptotic to $Hd_1(T)e^{-3T}/\pi$, although both the time interval and the
proportion tend to zero.
On each pole-free component of $I$, $g'(x)=-\sum_\gamma m_\gamma/(x-\gamma)^2<0$, so the occupation measure has
density $w_I(v)=\sum_{x\in I^\circ:g(x)=v}1/\abs{g'(x)}$, and the theorem is equivalent to
$\abs{w_I(v)-\int_Ip_{q,t}(v)\,\dd t}\le Cp^*_{q,T}(v)/\log d$ for almost every $v$, and at every $v$ different
from the finite values of $g$ at the endpoints: for $v=0$ a weighted count of critical points of $F$, for other
levels of zeros of $F'-vF$. With $E=(v,\infty)$ and endpoints not poles, $v\abs{\set{x\in I:g(x)>v}}\to\sum_{\gamma\in I}m_\gamma$
as $v\to\infty$, since $g(x)=m_\gamma/(x-\gamma)+O(1)$ at each pole, and multiplying by $v$ gives
\[
 \Bigl|\sum_{\gamma\in I}m_\gamma-\frac1\pi\int_Id_q(t)\,\dd t\Bigr|\ll\frac d{\log d},
\]
the order of the classical bound for the zero count under GRH, for which sharper constants are known
\cite{CarneiroChandeeMilinovich2015Note}; here the whole value line is treated in one measure inequality.

\begin{corollary}[Local Cauchy law on shrinking windows; \stG]\label{zer:thm:localcauchy}
Let $F$ be one of the functions \eqref{zer:eq:phasereal}, $g=F'/F$, $u<v$ fixed, and $U_T$ uniform on
$[uT,vT]$; put $Z_T=2g(U_T)/\log T$. If $\alpha_T<\beta_T$ are real with $\alpha_T,\beta_T\to c\in\R$, with
no lower bound on $\beta_T-\alpha_T$, then
\[
 \Pr\set{\alpha_T<Z_T<\beta_T}\sim\frac{\beta_T-\alpha_T}{\pi(1+c^2)} .
\]
Equivalently, if $a_T<b_T$, $\delta_T=b_T-a_T$ and $a_T/\log T,b_T/\log T\to d\in\R$, then
\[
 \abs{\set{t\in[uT,vT]:a_T<g(t)<b_T}}\sim\frac{2(v-u)}{\pi(1+4d^2)}\,\frac{\delta_TT}{\log T}.
\]
In particular, for every $\eps(T)>0$ with $\eps(T)=o(\log T)$,
\[
 \abs{\set{t\in[T,2T]:\abs{F'(t)/F(t)}<\eps(T)}}\sim\frac{4\eps(T)T}{\pi\log T},
\]
for example with $\eps(T)=e^{-T}$ and $F(t)=\xi(\tfrac12+it)$.
\end{corollary}
\begin{proof}
The first form follows from the second with $a_T=\alpha_T\log T/2$, $b_T=\beta_T\log T/2$, $d=c/2$, and the
last with $a_T=-\eps(T)$, $b_T=\eps(T)$, $d=0$. Put $E_T=(a_T,b_T)$, let $q$ be the conductor of $F$, and
let $N(J)$ be the measure of $\set{t\in J:g(t)\in E_T}$. Fix $t_0$ large. For $y\in E_T$ and $t\ge t_0$,
$p_{q,t}(y)\le1/(\pi d_q(t))$, and for fixed $0<\eps<w$, uniformly for $t\in[\eps T,wT]$ and $y\in E_T$,
$y/d_q(t)\to2d$, hence $p_{q,t}(y)=(1+o(1))\,2/(\pi(1+4d^2)\log T)$. Cover $[t_0,wT]$ by $O(\log T)$
dyadic intervals $I\subseteq[T',2T']$. On each, \cref{zer:thm:uniformcauchy} gives
$N(I)=\int_Ip_{q,t}(E_T)\,\dd t+O(\delta_T)$, since $p^*_{q,T'}(E_T)\le\delta_T/(\pi d_q(T'))$, and the model
integral over $[t_0,\eps T]$ is $O(\eps\delta_TT/\log T)$. Letting $\eps\to0$, for $0\le u<v$,
\[
 N\bigl([\max(uT,t_0),vT]\bigr)=\frac{2(v-u)}{\pi(1+4d^2)}\,\frac{\delta_TT}{\log T}\bigl(1+o(1)\bigr).
\]
Every error is proportional to $\delta_T$, so the rate of $\delta_T\to0$ does not enter. Negative times
reduce to positive ones: $F(-z)$ is the function \eqref{zer:eq:phasereal} of the conjugate character
($F_\zeta$ itself), of the same conductor and with logarithmic derivative $-g(-z)$, so $N([-vT,-uT])$ for $F$
is the corresponding measure for $F(-z)$ and the interval $(-b_T,-a_T)$, whose limit ratio $-d$ gives the
same constant. Finally $g$ has finitely many poles in $[-t_0,t_0]$, and on each pole-free component
$-g'(x)$, given by the sum after \eqref{zer:eq:logderim}, is at least $(t_0+\abs{\gamma_1})^{-2}$ for any fixed
zero $\gamma_1$ of $F$; so $N([-t_0,t_0])=O_{t_0}(\delta_T)$. Adding the pieces gives the second form for all
fixed $u<v$.
\end{proof}

A weak Cauchy limit gives no relative error for intervals whose length tends to zero, in particular
for a length $e^{-T}$; \cref{zer:thm:uniformcauchy} gives one, also on short time intervals. The
statements concern a continuous sampling variable and say nothing about what a discrete grid resolves.

\begin{theorem}[Convergent displacement remainder; \stG]\label{zer:thm:finitepart}
Fix the inputs and $A,B$ with $B\succeq0$, and $t_0$ so large that $D(t)\succ0$ for $t\ge t_0$. There is a real
constant $C(t_0)$ with
\[
 \int_{t_0}^T\sigma_{A,B}(t)\,\dd t=\int_{t_0}^T\rho(t)\,\dd t+C(t_0)+O\Bigl(\frac1{\log T\log\log T}\Bigr).
\]
For $B\ne0$ the improper integral defining $C(t_0)$ converges only conditionally.
\end{theorem}

A dyadic summation of \cref{zer:thm:sharpinterval} would give the divergent series $\sum_k1/(k\log k)$; a separate
contour argument makes the bound uniform from $a$ to every $b\ge a$.

\begin{proof}
Implied constants may depend on the data. Put $G_1(z)=G_0(z)+C_1/z$ with $C_1$ diagonal and real entries
$c_j=a_j/2-\tfrac14+2\cdot\mathbf 1\set{F_j=F_\zeta}$, $a_\zeta=0$: inserting $\psi(w)=\zerLog w-1/(2w)+O(\abs w^{-2})$ at
$w=(\tfrac12+a_j-iz)/2$, and $-i(1/s+1/(s-1))=2/z+O(\abs z^{-2})$ for the zeta polynomial, gives
$\norm{G(z)-G_1(z)}\ll2^{-y}+\abs z^{-2}$ uniformly for $x=\Re z$ large and $y=\Im z\ge1$.

Take $a$ large and $h=C_2/\log\log a$ with $C_2$ a large constant. For every $t\ge a$ the proof of
\cref{zer:lem:strip}, frozen at $T=t$, applies for $h\le y\le1$, and the resolvent calculation with scale
$\ell_t=\tfrac12\log t$ gives
$\abs{\Theta(G(t+iy))-\Theta(G_1(t+iy))}\ll e^{-cy\log\log t}/\log t+(\log t)^{-2}$, since $\Omega_t\ll1/\log t$ and $C_1/z=O(1/t)$.
For $y\ge1$ the imaginary parts of $G,G_1$ are at most $-c\log\abs z\,I$, so the phase difference is
$O((2^{-y}+\abs z^{-2})/(\log\abs z)^2)$. Integrating in $y$ costs $O(1/(\log t\log\log t))$ at each endpoint.
For finite $b\ge a$, move the horizontal integrals of $\Theta(G)-\Theta(G_1)$ at heights $h,2h$ to infinity:
the top edge tends to zero and the vertical edges have the bounds just proved, uniformly in $b$. The
clipping error \eqref{zer:eq:clip} is $O(h/\log a)$, and moving $\Theta(G_1)$ from height $h$ or $2h$ to the
boundary costs $O(h\int_a^b\dd t/(t(\log t)^2))=O(h/\log a)$, since $G_1'(t+iy)=O(1/t)$ and differentiating the
inverse supplies two factors $O(1/\log t)$. Finally both inverses $(A+uB+G_k(t))^{-1}$, $k=0,1$, equal
$iI/\ell_t+O(\ell_t^{-2})$ uniformly in $u$, so
$\Theta(G_1(t))-\Theta(G_0(t))=\tr(BC_1)/(t\ell_t^2)+O(1/(t\ell_t^3))$, whose leading term is real; its imaginary part
costs $O(1/(\log a)^2)$ after integration. Hence $\sup_{b\ge a}\abs{\int_a^b(\sigma_{A,B}-\rho)}\ll(\log a\log\log a)^{-1}$,
and the Cauchy criterion gives the theorem. For $B\ne0$, $\rho(t)\sim\tr B/(\pi\ell_t)$, so eventually
$0<\rho<\tfrac12$ with divergent integral, while $\sigma_{A,B}$ is integer valued, so $\abs{\sigma_{A,B}-\rho}\ge\rho$.
\end{proof}

\begin{corollary}[Coefficients on a sharp window; \stG]\label{zer:cor:sharpexp}
Fix the data, let $\mathsf C=\zerdiag(\kappa_i)$ as in \cref{zer:thm:transportexp}, $A_*=A-\tfrac\pi4I$ and
$U=\mathsf C+iA_*$. For every fixed $m\ge1$,
\[
 \pi\rho(t)=\sum_{k=1}^m\frac{(-1)^{k+1}}{k\ell_t^k}\Im\tr\bigl((U+iB)^k-U^k\bigr)+O_m(\ell_t^{-m-1}),
\]
and for $[a,b]\subseteq[T,2T]$ of length $H$, the integral of the sum divided by $\pi$ approximates
$\int_a^b\sigma_{A,B}$ with error $O_m\bigl(H/(\log T)^{m+1}+1/(\log T\log\log T)\bigr)$. The first three
numerators, signs and factors $1/k$ included, are
\[
 \tr B,\qquad-\tr(\mathsf CB),\qquad\tr(\mathsf C^2B)-\tr(A_*^2B)-\tr(A_*B^2)-\tfrac13\tr(B^3).
\]
\end{corollary}
\begin{proof}
$A+G_0(t)=-i\ell_t(I+U/\ell_t)$; expand the matrix logarithm near the identity, the scalar factors cancelling
in the difference, and combine with \cref{zer:thm:sharpinterval}; cyclicity of the trace suffices, with no
commutativity.
\end{proof}

So on $[T,2T]$ the window integrals have an expansion to every fixed order; it is the radial expansion of
\cref{zer:thm:transportexp} with $A$ replaced by $A-\tfrac\pi4I$, the real shift first visible in the third
coefficient.

For the zeta input with $A=0$, $B=1$ the old and new outputs are $F'$ and $F'+F$, and within each gap between
distinct zeros of $F$ the strictly decreasing $g$ runs from $+\infty$ to $-\infty$, so the contribution is the
distance from the solution of $g=0$ to that of $g=-1$, and
$\rho(t)=\pi^{-1}\bigl(\arctan\frac{1-\pi/4}{d_1(t)}-\arctan\frac{-\pi/4}{d_1(t)}\bigr)$. Grouping the first $23$ gaps,
with no fitted parameter:
\begin{center}\small
\begin{tabular}{rrrr}
\toprule
zero indices at the endpoints&measured length&$\int_I\rho$&difference\\\midrule
1 and 8&10.904081&10.451136&0.452945\\
8 and 16&6.415915&6.237236&0.178679\\
16 and 24&4.776198&4.726985&0.049213\\
\bottomrule
\end{tabular}
\end{center}
For instance the first gap runs from $14.13472514$ to $21.02203964$, with crossings $15.58570859$ and
$18.99408372$. These are floating-point illustrations at low height; exact rational checks of the kernel
identity and of the first three coefficients on two noncommuting examples are the third file of
\begin{center}\artifact{dirichlet-sharp-window-transport/checks/example.py},\\
\artifact{dirichlet-sharp-window-transport/checks/example.json},\\
\artifact{dirichlet-sharp-window-transport/checks/check_identities.py}.\end{center}

The scale $1/\log d$ is the limit of arguments based only on the Cayley-transform bound: for
$g_k(z)=\cot(kz/2)$, which maps the upper half-plane into the lower one with $(g_k+i)/(g_k-i)=e^{ikz}$ and
distance $O(e^{-k})$ from $-i$ at height one, the event $0<g_k<1$ fills an interval of length $\pi/(2k)$ in each
period $2\pi/k$, its Cauchy probability is $\tfrac14$, and taking that interval as the window gives discrepancy
$3\pi/(8k)$, of exact order $1/\log d$ with $d=e^k$. The optimal arithmetic scale is open (\cref{cpl:sec:open}).
The quantitative sharp-interval remainder, uniform after weighting by the size of the matrix phase, and the
finite part are new relative to the weak Cauchy limits and the Herglotz value-distribution theory
\cite{Sodin2018Critical,NajnudelNikeghbali2026Cauchy,Pearson1994Value,AizenmanWarzel2015Cauchy}, whose relation to
the matrix spectral shift is classical \cite{BehrndtMalamudNeidhardt2008Scattering}.

\subsection{Open questions}\label{cpl:sec:open}
{\raggedright\sources{graph_coupled_zeros_2026-11-26; grh_exact_extensions_2026-11-27; grh_path_rigidity_2026-11-27; grh_inverse_spectral_2026-11-27; coupled-dirichlet-inverse-spectrum; dirichlet-coupled-zeros-and-prime-variances; dirichlet-sharp-window-transport}}

\begin{enumerate}
\item \emph{Couplings on a prescribed topology.} \Cref{zer:thm:minimax} determines the optimal graph
for every $W_p$ when all edges are allowed. Which conductances are optimal when the support of the
graph is prescribed, for instance a path or a tree, or with edge caps or integral conductances (the
lower bounds of \cref{zer:thm:minimax} remain valid, but equality needs the displayed optimizer), and
which graph is optimal for a particular
arithmetic packet of zeros rather than for the worst case given the degrees?
\item \emph{The inverse problem.} Mixed parities, unknown lists, decoding and robustness are settled by
\cref{zer:thm:window,zer:cor:blindall,zer:thm:decode,zer:thm:stabledecoder}, and a sufficient window is
computable (\cref{zer:rem:windowsearch}). Can the first sufficient window of \cref{zer:thm:window} be bounded
explicitly by a useful function of the conductors, the matrix norm and the corruption budgets, together with
the stability constant of \cref{zer:thm:stabledecoder}? What height and precision suffice when every reported
ordinate is rounded? Can the edge weights alone be recovered in time polynomial in the number of vertices?
What is the strongest discrepancy constant obtainable in \eqref{zer:eq:witness} from the Landau coefficients
and the counts, over all primes and nonnegative trigonometric polynomials, and can it be determined uniformly
for conductors up to a bound? Classify the Hermitian matrices with equal principal minors, which the zeros
cannot separate.
\item \emph{Sharp transport.} Can the finite part $C(t_0)$ of \cref{zer:thm:finitepart} be computed to
certified accuracy for a concrete pair of couplings, with explicit constants in the tail? Does the endpoint
estimate of \cref{zer:thm:sharpinterval} hold for windows at the mean-spacing scale outside a controlled
exceptional set of heights, and what is the optimal arithmetic scale, given that arguments using only the
Cayley-transform bound stop at $1/\log d$?
\end{enumerate}
